\documentclass[11pt,a4paper,oneside]{book}

% Encoding and fonts
\usepackage[utf8]{inputenc}
\usepackage[T1]{fontenc}
\usepackage{lmodern}

% Page geometry
\usepackage[
  top=2.5cm,
  bottom=2.5cm,
  left=2.5cm,
  right=2.5cm,
  headheight=14pt
]{geometry}

% Language
\usepackage[english]{babel}

% Colors
\usepackage{xcolor}
\definecolor{darkblue}{HTML}{2C3E50}
\definecolor{codebg}{HTML}{F5F5F5}
\definecolor{codeframe}{HTML}{DDDDDD}
\definecolor{linkcolor}{HTML}{2C3E50}
\definecolor{chaptercolor}{HTML}{2C3E50}

% Code listings
\usepackage{listings}
\lstset{
  backgroundcolor=\color{codebg},
  frame=single,
  rulecolor=\color{codeframe},
  basicstyle=\ttfamily\small,
  breaklines=true,
  breakatwhitespace=false,
  tabsize=4,
  showstringspaces=false,
  numbers=none,
  xleftmargin=0.5em,
  xrightmargin=0.5em,
  aboveskip=1em,
  belowskip=1em,
  columns=fullflexible,
  keepspaces=true,
  literate={á}{{\'a}}1 {é}{{\'e}}1 {í}{{\'i}}1 {ó}{{\'o}}1 {ú}{{\'u}}1
}

% Headers and footers
\usepackage{fancyhdr}
\pagestyle{fancy}
\fancyhf{}
\fancyhead[L]{\small\itshape\leftmark}
\fancyhead[R]{\small\thepage}
\fancyfoot{}
\renewcommand{\headrulewidth}{0.4pt}

% Plain pages (chapter starts)
\fancypagestyle{plain}{
  \fancyhf{}
  \fancyfoot[C]{\small\thepage}
  \renewcommand{\headrulewidth}{0pt}
}

% Chapter and section styling
\usepackage{titlesec}

\titleformat{\chapter}[display]
  {\normalfont\huge\bfseries\color{chaptercolor}}
  {\chaptertitlename\ \thechapter}{20pt}{\Huge}
\titlespacing*{\chapter}{0pt}{-20pt}{30pt}

\titleformat{\section}
  {\normalfont\Large\bfseries\color{chaptercolor}}
  {\thesection}{1em}{}

\titleformat{\subsection}
  {\normalfont\large\bfseries\color{chaptercolor}}
  {\thesubsection}{1em}{}

% Links
\usepackage[
  colorlinks=true,
  linkcolor=linkcolor,
  urlcolor=linkcolor,
  citecolor=linkcolor,
  bookmarks=true,
  bookmarksnumbered=true,
  pdfstartview=FitH
]{hyperref}

% Tables
\usepackage{longtable}
\usepackage{booktabs}
\usepackage{array}

% Graphics
\usepackage{graphicx}

% Pandoc code highlighting compatibility
\usepackage{framed}
\usepackage{fancyvrb}
\newenvironment{Shaded}{\begin{snugshade}}{\end{snugshade}}
\definecolor{shadecolor}{HTML}{F5F5F5}
\DefineVerbatimEnvironment{Highlighting}{Verbatim}{
  commandchars=\\\{\},
  fontsize=\small,
  xleftmargin=0.5em
}
\newcommand{\KeywordTok}[1]{\textcolor[HTML]{007020}{\textbf{#1}}}
\newcommand{\DataTypeTok}[1]{\textcolor[HTML]{902000}{#1}}
\newcommand{\DecValTok}[1]{\textcolor[HTML]{40A070}{#1}}
\newcommand{\BaseNTok}[1]{\textcolor[HTML]{40A070}{#1}}
\newcommand{\FloatTok}[1]{\textcolor[HTML]{40A070}{#1}}
\newcommand{\ConstantTok}[1]{\textcolor[HTML]{880000}{#1}}
\newcommand{\CharTok}[1]{\textcolor[HTML]{4070A0}{#1}}
\newcommand{\SpecialCharTok}[1]{\textcolor[HTML]{4070A0}{#1}}
\newcommand{\StringTok}[1]{\textcolor[HTML]{4070A0}{#1}}
\newcommand{\VerbatimStringTok}[1]{\textcolor[HTML]{4070A0}{#1}}
\newcommand{\SpecialStringTok}[1]{\textcolor[HTML]{BB6688}{#1}}
\newcommand{\ImportTok}[1]{#1}
\newcommand{\CommentTok}[1]{\textcolor[HTML]{60A0B0}{\textit{#1}}}
\newcommand{\DocumentationTok}[1]{\textcolor[HTML]{BA2121}{\textit{#1}}}
\newcommand{\AnnotationTok}[1]{\textcolor[HTML]{60A0B0}{\textbf{\textit{#1}}}}
\newcommand{\CommentVarTok}[1]{\textcolor[HTML]{60A0B0}{\textbf{\textit{#1}}}}
\newcommand{\OtherTok}[1]{\textcolor[HTML]{007020}{#1}}
\newcommand{\FunctionTok}[1]{\textcolor[HTML]{06287E}{#1}}
\newcommand{\VariableTok}[1]{\textcolor[HTML]{19177C}{#1}}
\newcommand{\ControlFlowTok}[1]{\textcolor[HTML]{007020}{\textbf{#1}}}
\newcommand{\OperatorTok}[1]{\textcolor[HTML]{666666}{#1}}
\newcommand{\BuiltInTok}[1]{#1}
\newcommand{\ExtensionTok}[1]{#1}
\newcommand{\PreprocessorTok}[1]{\textcolor[HTML]{BC7A00}{#1}}
\newcommand{\AttributeTok}[1]{\textcolor[HTML]{7D9029}{#1}}
\newcommand{\RegionMarkerTok}[1]{#1}
\newcommand{\InformationTok}[1]{\textcolor[HTML]{60A0B0}{\textbf{\textit{#1}}}}
\newcommand{\WarningTok}[1]{\textcolor[HTML]{60A0B0}{\textbf{\textit{#1}}}}
\newcommand{\AlertTok}[1]{\textcolor[HTML]{FF0000}{\textbf{#1}}}
\newcommand{\ErrorTok}[1]{\textcolor[HTML]{FF0000}{\textbf{#1}}}
\newcommand{\NormalTok}[1]{#1}

% Tight list (pandoc compatibility)
\providecommand{\tightlist}{%
  \setlength{\itemsep}{0pt}\setlength{\parskip}{0pt}}

% Copyleft symbol (mirrored ©)
\providecommand{\textcopyleft}{\reflectbox{\copyright}}


% No paragraph indent, space between paragraphs
\setlength{\parindent}{0pt}
\setlength{\parskip}{0.6em}

% Quote environment styling (for epigraph)
\usepackage{etoolbox}
\AtBeginEnvironment{quote}{%
  \itshape\color{darkblue}%
}

% PDF metadata
\hypersetup{
  pdftitle={Ten Brief Lessons on Web Programming},
  pdfauthor={Giuseppe Costanzi},
  pdfsubject={Web Programming},
  pdfkeywords={PHP, JavaScript, HTML, CSS, SQLite, web development}
}

\begin{document}

% --- COVER PAGE ---
\begin{titlepage}
\newgeometry{margin=0pt}
\pagecolor{darkblue}
\begin{center}

\vspace*{6cm}

{\fontsize{28}{34}\selectfont\bfseries\color{white}
Ten Brief Lessons\\[0.4cm]on Web Programming}

\vspace{0.8cm}

{\large\color{white!80}
A concept notebook on web programming}

\vspace{2cm}

{\small\color{white!60}\itshape
``There are 10 types of people in the world:\\
those who understand binary, and those who don't.''}

\vspace{0.3cm}

{\small\color{white!80}
Which side are you on?}

\vspace{\fill}

{\large\color{white}
Giuseppe Costanzi}

\vspace{2cm}

\end{center}
\restoregeometry
\nopagecolor
\end{titlepage}

% --- BLANK PAGE ---
\thispagestyle{empty}
\vspace*{\fill}
\begin{center}
\textit{\color{gray}This page is intentionally left blank.}
\end{center}
\vspace*{\fill}
\newpage

% --- COPYRIGHT / COLOPHON ---
\thispagestyle{empty}
\vspace*{\fill}

{\small

\textbf{Ten Brief Lessons on Web Programming}\\
\textit{A concept notebook on web programming}

\vspace{1em}

Copyleft \textcopyleft\ 2026 Giuseppe Costanzi \textit{aka} \texttt{1966bc}

\vspace{1em}

Version 1.0 --- First edition, MMXXVI AD / MMDCCLXXIX \textit{ab urbe condita}

\vspace{1em}

This book is free software: you can redistribute it and/or modify
it under the terms of the GNU General Public License as published by
the Free Software Foundation, either version 3 of the License, or
(at your option) any later version.

\vspace{1em}

Source code: \url{https://github.com/1966bc/bibliotheca}

\vspace{1em}

Built with pure PHP, JavaScript, HTML, CSS and SQLite.\\
Typeset with \LaTeX\ via Pandoc.\\
No frameworks were harmed in the making of this book.

\vspace{1em}

\texttt{giuseppe.costanzi@gmail.com}

\vspace{1em}

{\footnotesize\itshape
Milky Way Galaxy, Solar System, third planet from the Sun\\
41\textdegree{}53'35''N \enspace 12\textdegree{}28'58''E --- Roma, Italia}
}

\newpage

% --- DEDICATION ---
\thispagestyle{empty}
\vspace*{5cm}
\begin{center}
\itshape\large
To the next generation of programmers:\\
learn to read, write and compute ---\\
before it's too late.
\end{center}
\newpage

% --- TABLE OF CONTENTS ---
{
% Add some stretch to help LaTeX distribute TOC entries better
\setlength{\parskip}{0.1em plus 0.1em}
\tableofcontents
}
\newpage

% --- BODY ---
% Start chapter numbering from 0
\setcounter{chapter}{-1}

\hypertarget{who-is-this-book-for}{%
\chapter*{Who Is This Book For}\label{who-is-this-book-for}}
\addcontentsline{toc}{chapter}{Who Is This Book For}

For the young programmer --- or the curious one at any age ---
who wants to understand how a web application actually works, in
anno Domini 2026. From the database to the browser, with nothing
in between but your own code.

You do not need experience with web development. But you should
be comfortable with basic programming: variables, functions,
loops, conditionals. You should know what a class is and how
objects work --- we use Object-Oriented Programming throughout,
and we will not stop to explain inheritance or encapsulation.

If you have never written a line of code, start there first.
Python is a good choice. Then come back.

\hypertarget{what-you-will-need}{%
\section*{What you will need}\label{what-you-will-need}}
\addcontentsline{toc}{section}{What you will need}

\begin{itemize}
\tightlist
\item
  A \textbf{Linux} machine (we use Debian, but any distribution
  works)
\item
  \textbf{Apache}, \textbf{PHP} and \textbf{SQLite} installed
\item
  A \textbf{text editor} (we recommend Sublime Text)
\item
  A \textbf{terminal} and a \textbf{web browser}
\item
  \textbf{Curiosity}
\end{itemize}

\hypertarget{what-you-will-learn}{%
\section*{What you will learn}\label{what-you-will-learn}}
\addcontentsline{toc}{section}{What you will learn}

\begin{itemize}
\tightlist
\item
  How HTTP works --- requests, responses, status codes
\item
  How to design a relational database with SQLite
\item
  How to structure an application with Model-View-Controller
\item
  How to build a REST API with pure PHP
\item
  How to build a dynamic frontend with pure JavaScript
\item
  How to implement CRUD operations from click to database row
\item
  How to validate, sanitize, and normalize user input
\item
  How to handle permissions, soft deletes, and dependencies
\item
  How to debug when things go wrong
\item
  How to secure your application: authentication, CSRF, headers,
  HTTPS
\end{itemize}

\hypertarget{what-this-book-is-not}{%
\section*{What this book is not}\label{what-this-book-is-not}}
\addcontentsline{toc}{section}{What this book is not}

This is not a reference manual. It is a journey. Each chapter
builds on the previous one. Read in order. Type the code
yourself. Break it. Fix it. That is how you learn.

We use no frameworks, no libraries, no build tools. Not because
they are bad --- but because you need to understand what they
replace before you use them.

\newpage

\hypertarget{prelude}{%
\chapter{Prelude}\label{prelude}}

\hypertarget{what-is-this}{%
\section{What is this}\label{what-is-this}}

The best way to learn something is by imitation. Study a real
application, understand how it works, then build your own.

Bibliotheca is a complete web application --- a book catalog.
Simple enough to hold in your head, rich enough to teach you the
fundamentals of web development.

The entire source code is on GitHub:
https://github.com/1966bc/bibliotheca

Read it. You will read more code than you write in your career
--- and reading good code is how you learn to write it. Reading
\emph{why} it is written that way is even better.

\hypertarget{who-is-this-for}{%
\section{Who is this for?}\label{who-is-this-for}}

For the young programmer's mind --- or the curious one at any age
--- open to the wonder of understanding how things actually work.
A web application, in anno Domini 2026. From the database to the
browser, with nothing in between but your own code.

You do not need experience with web development, but you should
know the basics of programming: variables, functions, loops,
conditionals. You should also have some intuition of how a
website works: that a browser asks a server for something, and
the server answers. And above all, the hacker's instinct: the
urge to open things up and see how they work.

\hypertarget{how-to-read-this}{%
\section{How to read this}\label{how-to-read-this}}

Read in order and try to absorb the key concepts. The lessons are
brief but there is a lot to take in. Look at the application as
you go --- piece by piece, the puzzle will come to life.

\hypertarget{what-you-will-need-1}{%
\section{What you will need}\label{what-you-will-need-1}}

\begin{itemize}
\tightlist
\item
  A \textbf{Linux} machine (we use Debian 12 Bookworm, but any
  distribution works)
\item
  A \textbf{LAMP} stack: Linux, Apache, MySQL/MariaDB, PHP ---
  even though we use SQLite instead of MySQL, Apache and PHP are
  essential. Install them with:
\end{itemize}

\begin{Shaded}
\begin{Highlighting}[]
\FunctionTok{sudo}\NormalTok{ apt install apache2 php libapache2{-}mod{-}php php{-}sqlite3}
\end{Highlighting}
\end{Shaded}

\begin{itemize}
\tightlist
\item
  A \textbf{text editor} --- we recommend
  \href{https://www.sublimetext.com/}{Sublime Text}, fast and
  lightweight. You do not need a heavy IDE to write good code.
\item
  A \textbf{terminal}
\item
  A \textbf{web browser}
\item
  \textbf{Curiosity}
\end{itemize}

One note about the terminal. Throughout this book, we will use
the command line for everything we can: creating files, running
queries, testing the API, checking permissions, managing git.
This is deliberate. The terminal is not a fallback for when the
GUI is missing. It is the primary tool of a programmer. Learn to
think in commands.

Let the show begin.

\newpage

\hypertarget{introduction}{%
\chapter{Introduction}\label{introduction}}

\hypertarget{the-shape-of-things-to-come}{%
\section{The Shape of Things to
Come}\label{the-shape-of-things-to-come}}

Projects begin in the mind, not on a computer. Then they take
shape with pen on paper. Because before writing a single line of
code, we need to define our rules and our domain. This is how
real projects are built.

As your thoughts take shape, write them down in your own words.
Reading and writing are two different cognitive processes.
Reading lets you follow someone else's reasoning. Writing forces
you to build your own. You will be surprised by the gap between
``I understood it'' and ``I can explain it.'' That gap is where
learning happens. Otherwise you end up like Augustine of Hippo:
\emph{``If no one asks me, I know; if I wish to explain it to one
who asks, I do not know.''}

\hypertarget{some-useful-principles-to-guide-our-decisions}{%
\section{Some useful principles to guide our
decisions}\label{some-useful-principles-to-guide-our-decisions}}

\begin{itemize}
\tightlist
\item
  \textbf{Böhm-Jacopini theorem} --- Sequence, selection
  (if/else), iteration (for/while). No magic, no tricks.
\item
  \textbf{KISS} --- Keep It Simple, Stupid. Always the simplest
  solution that works.
\item
  \textbf{DRY} --- Don't Repeat Yourself.
\item
  \textbf{YAGNI} --- You Aren't Gonna Need It. Don't build what
  you don't need yet.
\item
  \textbf{SoC} --- Separation of Concerns. Each part of the code
  has one job.
\item
  \textbf{Least Surprise} --- Code should behave as the reader
  expects.
\item
  \textbf{Fail Fast} --- If something goes wrong, stop
  immediately.
\item
  \textbf{Fail Safe} --- When you fail, fail securely.
\end{itemize}

\hypertarget{coding-standards}{%
\section{Coding standards}\label{coding-standards}}

We follow consistent conventions throughout the project:

\begin{itemize}
\tightlist
\item
  \textbf{Indentation}: 4 spaces everywhere (PHP, JavaScript,
  HTML, CSS, SQL). No tabs.
\item
  \textbf{Classes}: \texttt{PascalCase} --- \texttt{Book},
  \texttt{PublishersView}, \texttt{CategoryForm}.
\item
  \textbf{Methods and variables}: \texttt{camelCase} ---
  \texttt{getAll}, \texttt{findById}, \texttt{bookTitle}.
\item
  \textbf{Constants}: \texttt{UPPER\_SNAKE\_CASE} ---
  \texttt{DB\_PATH}, \texttt{API\_URL}.
\item
  \textbf{Files}: \texttt{snake\_case}, lowercase. Plural for
  collections (\texttt{publishers.php}), singular for single
  entity (\texttt{publisher.php}). PHP class files match the
  class name (\texttt{Book.php}).
\item
  \textbf{SQL}: keywords \texttt{UPPERCASE}, tables and columns
  \texttt{snake\_case}. Always prepared statements with named
  parameters (\texttt{:id}, \texttt{:title}).
\item
  \textbf{PHP}: \texttt{declare(strict\_types=1)} at the top of
  every file. Opening brace on the same line for control
  structures, next line for classes and methods.
\item
  \textbf{JavaScript}:
  \texttt{\textquotesingle{}use\ strict\textquotesingle{}} at the
  top. Always \texttt{const} or \texttt{let}, never \texttt{var}.
  Always \texttt{async/await}, never \texttt{.then()} chains.
  Semicolons required.
\item
  \textbf{HTML}: double quotes for attributes,
  \texttt{kebab-case} for IDs and classes, semantic tags.
\item
  \textbf{Comments}: English only. Explain \emph{why}, not
  \emph{how}.
\end{itemize}

\hypertarget{our-tools-the-stack}{%
\section{Our tools, the stack}\label{our-tools-the-stack}}

\begin{longtable}[]{@{}lll@{}}
\toprule
Layer & Technology & Notes \\
\midrule
\endhead
Backend & PHP (pure) & No frameworks \\
Frontend & JavaScript (pure) & No frameworks \\
Markup & HTML (pure) & No template engines \\
Style & CSS (pure) & No preprocessors \\
Database & SQLite & Via PDO \\
\bottomrule
\end{longtable}

Pure languages only. We use nothing but the native capabilities
of each language. No extra tools to learn, no abstractions to
cloud your understanding. Stay pure, stay focused.

\hypertarget{the-terminal}{%
\section{The terminal}\label{the-terminal}}

There is one tool that sits above all others: the command line.
Every operation in this book can be done from the terminal, and
most of them \emph{should} be. Creating a database, inserting
seed data, testing an API endpoint, checking file permissions,
committing code --- all from a text prompt.

Why? Because a graphical interface hides what is happening. A
command shows it. When you type
\texttt{chmod\ 664\ bibliotheca.db}, you know exactly what you
are changing. When you right-click and open a properties dialog,
you are trusting someone else's abstraction.

More importantly: a production server has no desktop, no file
manager, no mouse. It has a terminal and nothing else. If you can
only work with a GUI, you cannot work on a real server. The
programmers who are comfortable on the command line are the ones
who can diagnose a problem at 2 AM on a server they have never
seen before.

The command line has been the programmer's primary interface
since the DEC VT100 in 1978. Its screen had 80 columns and 24
rows --- and that is where the 80-character line convention comes
from, the same convention you still see in coding standards
today. Graphical desktops came and went. The terminal is still
here, because it works. In programming, traditions exist for a
reason.

Make the terminal your first choice, not your last resort.

Each of these technologies has extensive documentation online.
Study them separately. This book shows how the pieces fit
together. You will find all references in the Bibliography
appendix.

\hypertarget{the-body-of-a-web-application}{%
\section{The body of a web
application}\label{the-body-of-a-web-application}}

Let us begin with an abstraction. Before reading about
architecture, look at your own body.

\textbf{The skeleton (HTML)} --- Strip everything away: skin,
clothes, muscles. What remains is structure. A skull, a spine,
ribs. HTML is this: it says \emph{what exists} --- a heading, a
table, a form, a navigation menu. A website with no CSS and no
JavaScript still works. Try it: disable both in your browser. The
page is ugly, naked, but readable. That is the skeleton test.

\textbf{The skin (CSS)} --- Now dress the skeleton. Skin gives
visible shape to bones. Clothes add style, color, identity. CSS
does exactly this. But if you change clothes, the bones
underneath do not change.

\textbf{The nervous system (JavaScript)} --- Bones and skin alone
are a statue. Beautiful, but still. The nervous system carries
signals: from senses to brain, from brain to muscles. JavaScript
moves data: it takes the user's click on ``Save'', transmits it
to the server, receives the response, and updates the skeleton.
It does not decide \emph{what} to show or \emph{how} it looks. It
connects and transports.

\textbf{The internal organs (PHP + Database)} --- The heart
pumps, the liver filters, the kidneys purify. You cannot see
them, but without them the body dies. The PHP backend and the
SQLite database receive requests, validate data, save it,
retrieve it, protect its integrity. The user does not know they
exist --- and that is how it should be.

When layers mix --- HTML inside JavaScript, SQL inside HTML ---
it is like a body with the liver attached to the knee. Does it
work? Maybe. Is it healthy? Never.

\hypertarget{architecture}{%
\section{Architecture}\label{architecture}}

Model-View-Controller (MVC), simplified for clarity. Now you can
map the body metaphor to the technical terms:

\begin{itemize}
\tightlist
\item
  \textbf{Model} --- The organs. PHP classes that talk to the
  database and return data.
\item
  \textbf{Controller} --- The nervous system. PHP scripts (API)
  that receive requests and respond with JSON, plus JavaScript
  that carries data to the interface.
\item
  \textbf{View} --- The skeleton and skin. HTML structure + CSS
  appearance.
\end{itemize}

\newpage

\hypertarget{database}{%
\chapter{Database}\label{database}}

\hypertarget{why-start-here}{%
\section{Why start here?}\label{why-start-here}}

The database is the foundation. If the schema is solid,
everything else follows naturally:

Remember MVC?

From tables come Models. From Models come Controllers. From
Controllers come Views.

In the beginning was the table.

\hypertarget{sqlite}{%
\section{SQLite}\label{sqlite}}

We chose SQLite for its simplicity. No server to install, no
users to configure, no permissions to manage. One file and you
are ready.

PDO (PHP Data Objects) works the same way regardless of the
database engine. If one day you need to switch to MySQL or
PostgreSQL, you change the connection string and nothing else.

\hypertarget{the-domain}{%
\section{The domain}\label{the-domain}}

Bibliotheca is a book catalog. So the domain has four entities:

\begin{itemize}
\tightlist
\item
  \textbf{Publisher} --- who publishes the book (Einaudi, Penguin
  Books\ldots)
\item
  \textbf{Category} --- the classification (Fiction, Science,
  Philosophy\ldots)
\item
  \textbf{Author} --- who writes the book
\item
  \textbf{Book} --- the core entity
\end{itemize}

And their relationships:

\begin{itemize}
\tightlist
\item
  A book has \textbf{one} publisher (many-to-one)
\item
  A book has \textbf{one} category (many-to-one)
\item
  A book can have \textbf{many} authors, an author can have
  \textbf{many} books (many-to-many, via the
  \texttt{book\_author} junction table)
\end{itemize}

\hypertarget{schema-conventions}{%
\section{Schema conventions}\label{schema-conventions}}

Every table follows the same structure, a convention we impose on
ourselves:

\begin{enumerate}
\def\labelenumi{\arabic{enumi}.}
\tightlist
\item
  First field: primary key (\texttt{table\_id})
\item
  Second field(s): foreign keys (if any)
\item
  Data fields
\item
  Last field: \texttt{status} (soft delete --- 1 active, 0
  inactive)
\end{enumerate}

Take note: we always access data \textbf{by column name}, never
by position. Say it out loud. You may not fully understand why
yet, but one day you will be grateful for this choice.

\hypertarget{the-schema}{%
\section{The schema}\label{the-schema}}

\begin{Shaded}
\begin{Highlighting}[]
\NormalTok{publisher (publisher\_id, name, status)}
\KeywordTok{category}\NormalTok{  (category\_id, name, status)}
\NormalTok{author    (author\_id, first\_name, last\_name, birthdate, status)}
\NormalTok{book      (book\_id, publisher\_id, category\_id, title, pages, published, status)}
\NormalTok{book\_author (book\_id, author\_id)}
\end{Highlighting}
\end{Shaded}

Five tables. Four entities, one junction table. Simple enough to
hold in your head while reading the code.

\hypertarget{data-types}{%
\section{Data types}\label{data-types}}

The schema intentionally covers the fundamental data types:

\begin{itemize}
\tightlist
\item
  \textbf{TEXT} --- strings (title, name, first\_name,
  last\_name)
\item
  \textbf{INTEGER} --- numbers (pages, status, primary and
  foreign keys)
\item
  \textbf{TEXT as DATE} --- dates in ISO 8601 format (published,
  birthdate)
\end{itemize}

\hypertarget{setting-up-the-console}{%
\section{Setting up the console}\label{setting-up-the-console}}

SQLite has a command-line interface. A programmer should always
prefer the command line. It builds character, and keeps you away
from the dark side of the force. We configure it with a
\texttt{setconsole} file that sets headers, column mode, and
enables foreign keys:

\begin{verbatim}
sqlite3 bibliotheca.db -init sql/setconsole
\end{verbatim}

\hypertarget{file-structure}{%
\section{File structure}\label{file-structure}}

\begin{verbatim}
sql/
    setconsole          — SQLite console configuration
    run.sql             — scratch pad for quick queries
    ddl/
        create_table.sql  — the schema (CREATE TABLE statements)
    dml/
        insert.sql        — sample data (INSERT statements)
    dql/
        (queries will go here)
    bibliotheca.db      — the database file
\end{verbatim}

\begin{itemize}
\tightlist
\item
  \textbf{ddl/} --- Data Definition Language (structure)
\item
  \textbf{dml/} --- Data Manipulation Language (data)
\item
  \textbf{dql/} --- Data Query Language (queries)
\end{itemize}

Take your time in this folder. Read the files, chew on them. It
will pay off.

\newpage

\hypertarget{project-structure}{%
\chapter{Project Structure}\label{project-structure}}

\hypertarget{security-by-design}{%
\section{Security by design}\label{security-by-design}}

The most important decision in a web application structure is
what the browser can and cannot reach. Everything inside the
document root is accessible via URL. Everything outside is not.

This is not a feature. It is a security rule.

\hypertarget{the-directory-tree}{%
\section{The directory tree}\label{the-directory-tree}}

\begin{verbatim}
bibliotheca/
    docs/                       — the book you are reading
    sql/                        — database (schema, data, queries)
        ddl/                    — CREATE TABLE statements
        dml/                    — INSERT statements
        dql/                    — SELECT queries
        setconsole              — SQLite console configuration
        run.sql                 — scratch pad for quick queries
        bibliotheca.db          — the database file
    src/                        — PHP classes — NOT accessible from the web
        DBMS.php                — database wrapper
        Publisher.php           — Publisher model
        Category.php            — Category model
        Author.php              — Author model
        Book.php                — Book model
    public/                     — document root — ONLY folder exposed to the browser
        .htaccess               — URL rewriting rules
        index.php               — the router
        css/
            style.css
        js/
            publishers.js       — list view (plural)
            publisher.js        — form view (singular)
            ...                 — same pattern for each entity
        api/
            publishers.php      — JSON endpoint
            ...                 — one per entity
        pages/
            home.php
            publishers.php      — list page (plural)
            publisher.php       — form page (singular)
            ...                 — same pattern for each entity
            404.php
\end{verbatim}

\hypertarget{two-worlds}{%
\section{Two worlds}\label{two-worlds}}

\begin{itemize}
\tightlist
\item
  \textbf{Public} (\texttt{public/}) --- what the browser can
  reach. HTML, CSS, JavaScript, and the API endpoints.
\item
  \textbf{Private} (\texttt{src/}, \texttt{sql/}) --- what only
  PHP can reach.
\end{itemize}

\hypertarget{the-stargate}{%
\section{The Stargate}\label{the-stargate}}

There is a principle in programming: don't talk to strangers. The
Stargate is where we enforce it. Everything outside
\texttt{public/} is invisible to the browser. The gate decides
who passes and where they go.

When the browser requests \texttt{/publishers}, Apache looks for
a file or directory called \texttt{publishers} inside
\texttt{public/}. It does not exist. Before giving up with a 404,
Apache checks if there is an \texttt{.htaccess} file in that
directory with instructions. It finds ours, reads it, and follows
the rewrite rule written in it:

\begin{Shaded}
\begin{Highlighting}[]
\ExtensionTok{RewriteEngine}\CharTok{ }\KeywordTok{On}
\NormalTok{RewriteCond}\StringTok{ \%\{REQUEST\_FILENAME\} !{-}f}
\NormalTok{RewriteCond}\StringTok{ \%\{REQUEST\_FILENAME\} !{-}d}
\NormalTok{RewriteRule}\StringTok{ \^{}(.*)$ index.php?route=$1 [L,QSA]}
\end{Highlighting}
\end{Shaded}

Line by line:

\begin{enumerate}
\def\labelenumi{\arabic{enumi}.}
\tightlist
\item
  \textbf{RewriteEngine On} --- enable URL rewriting.
\item
  \textbf{RewriteCond !-f} --- only rewrite if the request is NOT
  an existing \textbf{f}ile (so \texttt{style.css} and
  \texttt{book.js} are served normally). The \texttt{-f} stands
  for \emph{file}.
\item
  \textbf{RewriteCond !-d} --- only rewrite if the request is NOT
  an existing \textbf{d}irectory. The \texttt{-d} stands for
  \emph{directory}.
\item
  \textbf{RewriteRule} --- capture the entire URL path and pass
  it to \texttt{index.php} as the \texttt{route} parameter.
  \texttt{{[}L{]}} means stop processing rules.
  \texttt{{[}QSA{]}} means append any existing query string (so
  \texttt{/book?id=3} becomes
  \texttt{index.php?route=book\&id=3}).
\end{enumerate}

The result: the user sees clean URLs (\texttt{/publishers},
\texttt{/book?id=3}), but every request actually goes through
\texttt{index.php}. This technique is called \textbf{routing}:
one entry point, many destinations. If a route is not in the
whitelist, it does not pass. It is the foundation of all modern
web applications.

Note: \texttt{.htaccess} requires Apache's \texttt{mod\_rewrite}
module, which is the engine that makes URL rewriting possible. If
clean URLs do not work, enable it with
\texttt{sudo\ a2enmod\ rewrite} and restart Apache.

\hypertarget{the-journey-of-a-request}{%
\section{The journey of a
request}\label{the-journey-of-a-request}}

\begin{verbatim}
Browser                          Server
──────                           ──────
URL: /publishers
  └── .htaccess rewrites to index.php
        └── loads pages/publishers.php
              └── loads js/publishers.js
                    │
                    │  fetch('/api/publishers.php')
                    │
                    └──────────────────► api/publishers.php (Controller)
                                            │
                                            │  require '../../src/Publisher.php'
                                            │
                                            └──► src/Publisher.php (Model)
                                                    │
                                                    │  PDO query
                                                    │
                                                    └──► sql/bibliotheca.db
\end{verbatim}

\newpage

\hypertarget{backend}{%
\chapter{Backend}\label{backend}}

In our architecture, PHP has one job: talk to the database and
return data. It lives outside \texttt{public/}, invisible to the
browser. It receives requests, processes them, and answers with
JSON. The internal organs, remember?

Most PHP books and tutorials embed PHP directly inside HTML. Do
not let that mislead you. PHP's real job is the backend:
validate, query, respond. For everything the user sees and
interacts with, there is JavaScript. This decoupling is
fundamental. The sooner you internalize it, the better.

\hypertarget{prerequisites}{%
\section{Prerequisites}\label{prerequisites}}

PHP needs the SQLite driver. Check if it is installed:

\begin{Shaded}
\begin{Highlighting}[]
\ExtensionTok{php} \AttributeTok{{-}m} \KeywordTok{|} \FunctionTok{grep} \AttributeTok{{-}i}\NormalTok{ sqlite}
\end{Highlighting}
\end{Shaded}

If empty, install it and restart Apache:

\begin{Shaded}
\begin{Highlighting}[]
\FunctionTok{sudo}\NormalTok{ apt install php{-}sqlite3}
\FunctionTok{sudo}\NormalTok{ systemctl restart apache2}
\end{Highlighting}
\end{Shaded}

\hypertarget{dbms-the-database-class}{%
\section{DBMS --- the database
class}\label{dbms-the-database-class}}

Open \texttt{src/DBMS.php} and read it. This is the only class
that talks to the database. Every query, every insert, every
transaction passes through it. It is the first file you should
study.

Most tutorials have you open a connection, write the query, and
fetch results on every single page. That is not how it works in
real projects. We write it once, in one place. Everyone else uses
it. That is DRY in action.

Naturally, it is a class we wrote ourselves. A thin wrapper
around PDO. It exposes two ways to execute SQL, matching what PDO
itself offers:

\begin{itemize}
\tightlist
\item
  \textbf{\texttt{query()}} --- for fixed SQL with no parameters.
  Example: \texttt{SELECT\ *\ FROM\ publisher\ ORDER\ BY\ name}.
\item
  \textbf{\texttt{fetchOne()} / \texttt{fetchAll()}} --- for SQL
  with named parameters, using prepared statements. Example:
  \texttt{WHERE\ publisher\_id\ =\ :id}.
\end{itemize}

And convenience methods for write operations (all via prepare):
\texttt{insert()}, \texttt{update()}, \texttt{delete()}.

It also provides \texttt{exec()} for statements that return no
rows (DDL, PRAGMA), and \texttt{beginTransaction()} /
\texttt{commit()} / \texttt{rollBack()} for atomic operations ---
when two or more statements must succeed together or fail
together.

As a safety net, \texttt{query()} and \texttt{exec()} reject SQL
that contains named parameters (\texttt{:name}). If you need
parameters, use prepared statements --- no exceptions.

Key choices in the constructor:

\begin{itemize}
\tightlist
\item
  \texttt{ERRMODE\_EXCEPTION} --- fail fast on errors.
\item
  \texttt{FETCH\_ASSOC} --- access data by column name, never by
  position.
\item
  \texttt{PRAGMA\ foreign\_keys\ =\ ON} --- SQLite ignores
  foreign keys by default.
\end{itemize}

\hypertarget{php-in-the-mvc}{%
\section{PHP in the MVC}\label{php-in-the-mvc}}

Remember MVC from Chapter 00? PHP lives in the Model
(\texttt{src/}) and in the Controller (\texttt{public/api/}).
JavaScript and HTML/CSS handle the rest.

\hypertarget{the-model}{%
\clearpage
\section{The Model}\label{the-model}}

Open \texttt{src/Publisher.php}. The Model receives the DBMS in
its constructor (\textbf{dependency injection}) and provides
methods to query one entity. It does not know where the database
lives or how it was opened. It just uses it.

The \texttt{status} column (1 = active, 0 = disabled) controls
visibility. \texttt{getAll} returns everything;
\texttt{getActive} filters by \texttt{status\ =\ 1}. Disabled
records appear greyed out in the list. Deletion is permanent
(\texttt{DELETE\ FROM}).

\hypertarget{the-controller}{%
\section{The Controller}\label{the-controller}}

Open \texttt{public/api/publishers.php}.

The Controller is the only PHP file the browser can reach. It
creates the DBMS and the Model, reads the HTTP request, calls the
Model, and returns JSON. Nothing else.

\hypertarget{testing-with-curl}{%
\section{Testing with curl}\label{testing-with-curl}}

PHP runs on the server. Unlike JavaScript, it does not need a
browser. Keep these two worlds separate in your mind. As you
grow, you will notice subtle exceptions, but for now: PHP is the
server, JavaScript is the browser. We can prove it right now from
the command line:

\begin{Shaded}
\begin{Highlighting}[]
\ExtensionTok{curl}\NormalTok{ http://localhost/bibliotheca/public/api/publishers.php}
\end{Highlighting}
\end{Shaded}

\begin{Shaded}
\begin{Highlighting}[]
\OtherTok{[}\FunctionTok{\{}\DataTypeTok{"publisher\_id"}\FunctionTok{:}\DecValTok{7}\FunctionTok{,}\DataTypeTok{"name"}\FunctionTok{:}\StringTok{"Addison{-}Wesley"}\FunctionTok{,}\DataTypeTok{"status"}\FunctionTok{:}\DecValTok{1}\FunctionTok{\}}\OtherTok{,}
 \FunctionTok{\{}\DataTypeTok{"publisher\_id"}\FunctionTok{:}\DecValTok{1}\FunctionTok{,}\DataTypeTok{"name"}\FunctionTok{:}\StringTok{"Adelphi"}\FunctionTok{,}\DataTypeTok{"status"}\FunctionTok{:}\DecValTok{1}\FunctionTok{\}}\OtherTok{,}
 \FunctionTok{\{}\DataTypeTok{"publisher\_id"}\FunctionTok{:}\DecValTok{2}\FunctionTok{,}\DataTypeTok{"name"}\FunctionTok{:}\StringTok{"Einaudi"}\FunctionTok{,}\DataTypeTok{"status"}\FunctionTok{:}\DecValTok{1}\FunctionTok{\}}\OtherTok{,}
 \FunctionTok{\{}\DataTypeTok{"publisher\_id"}\FunctionTok{:}\DecValTok{3}\FunctionTok{,}\DataTypeTok{"name"}\FunctionTok{:}\StringTok{"Feltrinelli"}\FunctionTok{,}\DataTypeTok{"status"}\FunctionTok{:}\DecValTok{1}\FunctionTok{\}}\OtherTok{,}
 \FunctionTok{\{}\DataTypeTok{"publisher\_id"}\FunctionTok{:}\DecValTok{8}\FunctionTok{,}\DataTypeTok{"name"}\FunctionTok{:}\StringTok{"Gallimard"}\FunctionTok{,}\DataTypeTok{"status"}\FunctionTok{:}\DecValTok{1}\FunctionTok{\}}\OtherTok{,}
 \FunctionTok{\{}\DataTypeTok{"publisher\_id"}\FunctionTok{:}\DecValTok{5}\FunctionTok{,}\DataTypeTok{"name"}\FunctionTok{:}\StringTok{"Laterza"}\FunctionTok{,}\DataTypeTok{"status"}\FunctionTok{:}\DecValTok{1}\FunctionTok{\}}\OtherTok{,}
 \FunctionTok{\{}\DataTypeTok{"publisher\_id"}\FunctionTok{:}\DecValTok{4}\FunctionTok{,}\DataTypeTok{"name"}\FunctionTok{:}\StringTok{"Mondadori"}\FunctionTok{,}\DataTypeTok{"status"}\FunctionTok{:}\DecValTok{1}\FunctionTok{\}}\OtherTok{,}
 \FunctionTok{\{}\DataTypeTok{"publisher\_id"}\FunctionTok{:}\DecValTok{6}\FunctionTok{,}\DataTypeTok{"name"}\FunctionTok{:}\StringTok{"Penguin Books"}\FunctionTok{,}\DataTypeTok{"status"}\FunctionTok{:}\DecValTok{1}\FunctionTok{\}}\OtherTok{]}
\end{Highlighting}
\end{Shaded}

Eight publishers, ordered by name. The backend works.

\newpage

\hypertarget{frontend}{%
\chapter{Frontend}\label{frontend}}

\hypertarget{the-flow}{%
\section{The flow}\label{the-flow}}

\begin{enumerate}
\def\labelenumi{\arabic{enumi}.}
\tightlist
\item
  The browser requests a URL (e.g., \texttt{/publishers}).
\item
  Apache rewrites it through \texttt{.htaccess} to
  \texttt{index.php}.
\item
  The router loads the right page from \texttt{pages/}.
\item
  The page includes a JavaScript file.
\item
  JavaScript calls \texttt{fetch} to the API, gets JSON back.
\item
  JavaScript builds the HTML from the data.
\end{enumerate}

The HTML page loads once. After that, all data exchange happens
through \texttt{fetch} calls that return JSON. The browser never
asks the server for a new page.

\hypertarget{routing}{%
\section{Routing}\label{routing}}

All URLs pass through a single entry point: \texttt{index.php}.
Apache makes this possible with one rewrite rule in
\texttt{.htaccess}:

\begin{verbatim}
RewriteRule ^(.+)$ index.php?route=$1 [QSA,L]
\end{verbatim}

Every request that does not match an existing file gets
rewritten. The URL \texttt{/publishers} becomes
\texttt{index.php?route=publishers}.

Inside \texttt{index.php}, the route is checked against a
whitelist:

\begin{Shaded}
\begin{Highlighting}[]
\VariableTok{$routes} \OperatorTok{=}\NormalTok{ [}\StringTok{\textquotesingle{}home\textquotesingle{}}\OtherTok{,} \StringTok{\textquotesingle{}login\textquotesingle{}}\OtherTok{,} \StringTok{\textquotesingle{}publishers\textquotesingle{}}\OtherTok{,} \StringTok{\textquotesingle{}publisher\textquotesingle{}}\OtherTok{,}
           \StringTok{\textquotesingle{}categories\textquotesingle{}}\OtherTok{,} \StringTok{\textquotesingle{}category\textquotesingle{}}\OtherTok{,} \StringTok{\textquotesingle{}authors\textquotesingle{}}\OtherTok{,} \StringTok{\textquotesingle{}author\textquotesingle{}}\OtherTok{,}
           \StringTok{\textquotesingle{}book\textquotesingle{}}\OtherTok{,} \StringTok{\textquotesingle{}about\textquotesingle{}}\NormalTok{]}\OtherTok{;}
\end{Highlighting}
\end{Shaded}

If the route is in the list and the file exists, the
corresponding page is loaded. Otherwise, the user gets a 404. No
user input ever becomes a file path directly.

\hypertarget{pages-plural-and-singular}{%
\section{Pages: plural and
singular}\label{pages-plural-and-singular}}

In our project, each entity has two pages, by convention:

\begin{itemize}
\tightlist
\item
  \textbf{Plural} (\texttt{publishers.php}): the list, with Edit
  and Delete buttons shown if the user is logged in.
\item
  \textbf{Singular} (\texttt{publisher.php}): the form, for
  adding or editing.
\end{itemize}

One page, one job. Separation of Concerns.

\hypertarget{fetch-the-bridge}{%
\section{fetch: the bridge}\label{fetch-the-bridge}}

\texttt{fetch} is the browser's native API for HTTP requests. It
replaced \texttt{XMLHttpRequest}, the old AJAX interface that
powered the first generation of dynamic web applications. The
idea is the same: JavaScript talks to the server without
reloading the page. The API is simpler.

\texttt{fetch} is natively asynchronous: the browser sends the
request and continues executing code without waiting for the
response. This is necessary because network calls are slow
compared to everything else, and blocking the browser would
freeze the page.

The \texttt{async/await} syntax lets us write asynchronous code
that reads like sequential code. Without it, we would need
callbacks or promise chains. With it, each line waits for the
previous one to finish, exactly as you would expect.

Reading data:

\begin{Shaded}
\begin{Highlighting}[]
\KeywordTok{async} \FunctionTok{load}\NormalTok{() \{}
    \KeywordTok{const}\NormalTok{ response }\OperatorTok{=} \ControlFlowTok{await} \FunctionTok{fetch}\NormalTok{(}\StringTok{\textquotesingle{}/api/publishers.php\textquotesingle{}}\NormalTok{)}\OperatorTok{;}
    \KeywordTok{const}\NormalTok{ publishers }\OperatorTok{=} \ControlFlowTok{await}\NormalTok{ response}\OperatorTok{.}\FunctionTok{json}\NormalTok{()}\OperatorTok{;}
\NormalTok{\}}
\end{Highlighting}
\end{Shaded}

Three lines. Request, parse, done. No callbacks, no libraries.
The \texttt{async} keyword on the method is required:
\texttt{await} only works inside an \texttt{async} function.

Sending data requires a bit more: the HTTP method, headers, and a
JSON body.

\begin{Shaded}
\begin{Highlighting}[]
\KeywordTok{const}\NormalTok{ payload }\OperatorTok{=}\NormalTok{ \{ }\DataTypeTok{name}\OperatorTok{:} \StringTok{\textquotesingle{}Adelphi\textquotesingle{}}\NormalTok{ \}}\OperatorTok{;}

\KeywordTok{const}\NormalTok{ response }\OperatorTok{=} \ControlFlowTok{await} \FunctionTok{fetch}\NormalTok{(}\StringTok{\textquotesingle{}/api/publishers.php\textquotesingle{}}\OperatorTok{,}\NormalTok{ \{}
    \DataTypeTok{method}\OperatorTok{:} \StringTok{\textquotesingle{}POST\textquotesingle{}}\OperatorTok{,}
    \DataTypeTok{headers}\OperatorTok{:}\NormalTok{ \{}
        \StringTok{\textquotesingle{}Content{-}Type\textquotesingle{}}\OperatorTok{:} \StringTok{\textquotesingle{}application/json\textquotesingle{}}\OperatorTok{,}
        \StringTok{\textquotesingle{}X{-}CSRF{-}Token\textquotesingle{}}\OperatorTok{:}\NormalTok{ token}
\NormalTok{    \}}\OperatorTok{,}
    \DataTypeTok{body}\OperatorTok{:} \BuiltInTok{JSON}\OperatorTok{.}\FunctionTok{stringify}\NormalTok{(payload)}
\NormalTok{\})}\OperatorTok{;}
\end{Highlighting}
\end{Shaded}

The payload is the data we want to send. \texttt{JSON.stringify}
converts the JavaScript object into a JSON string, because HTTP
only carries text. The CSRF token travels in the header. The data
travels in the body. The server reads both before doing anything.

\hypertarget{dom-building-the-page}{%
\section{DOM: building the page}\label{dom-building-the-page}}

The DOM (Document Object Model) is the browser's internal
representation of the HTML. Every tag becomes a node, every
attribute a property. JavaScript does not modify the HTML file.
It modifies this tree, and the browser redraws the screen.

Take this HTML from our publishers page:

\begin{Shaded}
\begin{Highlighting}[]
\KeywordTok{\textless{}table} \ErrorTok{id}\OtherTok{=}\StringTok{"publisher{-}table"}\KeywordTok{\textgreater{}}
    \KeywordTok{\textless{}thead\textgreater{}}
        \KeywordTok{\textless{}tr\textgreater{}}
            \KeywordTok{\textless{}th\textgreater{}}\NormalTok{Name}\KeywordTok{\textless{}/th\textgreater{}}
            \KeywordTok{\textless{}th\textgreater{}\textless{}/th\textgreater{}}
        \KeywordTok{\textless{}/tr\textgreater{}}
    \KeywordTok{\textless{}/thead\textgreater{}}
    \KeywordTok{\textless{}tbody\textgreater{}\textless{}/tbody\textgreater{}}
\KeywordTok{\textless{}/table\textgreater{}}
\end{Highlighting}
\end{Shaded}

The browser reads it and builds the DOM, a tree of objects:

\begin{verbatim}
table #publisher-table
├── thead
│   └── tr
│       ├── th  "Name"
│       └── th
└── tbody  (empty)
\end{verbatim}

JavaScript actually sees this tree, not the HTML text.

This is a crucial point. The \texttt{id} on the table and the
empty \texttt{tbody} are placeholders. They tell JavaScript
\emph{where} to inject the data retrieved by \texttt{fetch}. The
HTML defines the structure; JavaScript fills it. When we call
\texttt{document.querySelector(\textquotesingle{}\#publisher-table\ tbody\textquotesingle{})},
we get the \texttt{tbody} node, ready to receive rows.

A concrete example from our code. When the book form loads, it
populates a dropdown with publishers fetched from the API:

\begin{Shaded}
\begin{Highlighting}[]
\ControlFlowTok{for}\NormalTok{ (}\KeywordTok{let}\NormalTok{ i }\OperatorTok{=} \DecValTok{0}\OperatorTok{;}\NormalTok{ i }\OperatorTok{\textless{}}\NormalTok{ publishers}\OperatorTok{.}\AttributeTok{length}\OperatorTok{;}\NormalTok{ i}\OperatorTok{++}\NormalTok{) \{}
    \KeywordTok{const}\NormalTok{ option }\OperatorTok{=} \BuiltInTok{document}\OperatorTok{.}\FunctionTok{createElement}\NormalTok{(}\StringTok{\textquotesingle{}option\textquotesingle{}}\NormalTok{)}\OperatorTok{;}
\NormalTok{    option}\OperatorTok{.}\AttributeTok{value} \OperatorTok{=}\NormalTok{ publishers[i]}\OperatorTok{.}\AttributeTok{publisher\_id}\OperatorTok{;}
\NormalTok{    option}\OperatorTok{.}\AttributeTok{textContent} \OperatorTok{=}\NormalTok{ publishers[i]}\OperatorTok{.}\AttributeTok{name}\OperatorTok{;}
    \KeywordTok{this}\OperatorTok{.}\AttributeTok{selectPublisher}\OperatorTok{.}\FunctionTok{appendChild}\NormalTok{(option)}\OperatorTok{;}
\NormalTok{\}}
\end{Highlighting}
\end{Shaded}

This is the core mechanism: JavaScript fetches the data, creates
new HTML elements on the fly, places them in the right spot in
the DOM, and can style them through classes or attributes. All of
this happens without reloading the page. \texttt{createElement}
makes the node. \texttt{textContent} sets its text.
\texttt{appendChild} attaches it to the page. Three operations,
always in this order.

Remember, never \texttt{innerHTML} with external data. That is an
XSS vulnerability. \texttt{textContent} is safe by design: it
treats everything as plain text, never as HTML.

\hypertarget{the-javascript-classes}{%
\section{The JavaScript classes}\label{the-javascript-classes}}

In our project each HTML page has a dedicated JavaScript class.
Every class has a \textbf{constructor} that stores references to
DOM elements and binds event listeners. The methods that follow
depend on the type of page.

\hypertarget{table-pages}{%
\subsection{Table pages}\label{table-pages}}

The table pages (\texttt{Publishers}, \texttt{Categories},
\texttt{Authors}, \texttt{Books}) display all records. These have
two methods:

\begin{itemize}
\tightlist
\item
  \textbf{\texttt{load()}}: fetches all records from the API.
\item
  \textbf{\texttt{render(data)}}: builds the table rows from the
  data.
\end{itemize}

\hypertarget{form-pages}{%
\subsection{Form pages}\label{form-pages}}

The form pages (\texttt{Publisher}, \texttt{Category},
\texttt{Author}, \texttt{Book}) handle create, edit, and delete.
These have more methods:

\begin{itemize}
\tightlist
\item
  \textbf{\texttt{checkEdit()}}: reads the \texttt{?id=}
  parameter from the URL. If present, the form is in edit mode
  and calls \texttt{load(id)}. If absent, the form is in create
  mode and stays empty.
\item
  \textbf{\texttt{load(id)}}: fetches one record from the API.
\item
  \textbf{\texttt{render(data)}}: populates the input fields with
  the fetched data.
\item
  \textbf{\texttt{save()}}: validates the input, builds a
  \texttt{payload} object, then calls \texttt{fetch} with POST
  (create) or PUT (update).
\item
  \textbf{\texttt{remove()}}: asks for confirmation, then sends a
  DELETE request.
\end{itemize}

Notice that \texttt{save} and \texttt{remove} use the same
\texttt{fetch} call structure. The only thing that changes is the
HTTP method:

\begin{longtable}[]{@{}ll@{}}
\toprule
Operation & Method \\
\midrule
\endhead
Create & \texttt{POST} \\
Update & \texttt{PUT} \\
Delete & \texttt{DELETE} \\
\bottomrule
\end{longtable}

Every HTTP request carries a \textbf{method}: a word that tells
the server what the client wants to do. Most introductory books
only teach two: GET to read a page, POST to send a form. But HTTP
defines more, and each has a precise role:

\begin{itemize}
\tightlist
\item
  \textbf{GET}: retrieve data. The browser sends this when you
  type a URL or click a link.
\item
  \textbf{POST}: create a new resource. The server receives data
  it has never seen before.
\item
  \textbf{PUT}: replace an existing resource. The server updates
  the record that matches the given ID.
\item
  \textbf{DELETE}: remove a resource. The server deletes the
  record.
\end{itemize}

There are others (PATCH, HEAD, OPTIONS), but these four are all
you need for a CRUD application. In the next chapter we will see
how each method maps to a CRUD operation.

The URL is always the same (\texttt{this.API}). The server reads
the method to decide what to do. This is \textbf{REST}:
Representational State Transfer. The name is academic, but the
idea is simple. Remember the constructor of our classes?

\begin{Shaded}
\begin{Highlighting}[]
\KeywordTok{this}\OperatorTok{.}\AttributeTok{API} \OperatorTok{=} \StringTok{\textquotesingle{}/bibliotheca/public/api/publishers.php\textquotesingle{}}\OperatorTok{;}
\end{Highlighting}
\end{Shaded}

That URL is what REST calls a \textbf{resource}. Every time we
call \texttt{fetch}, we point to that same URL:

\begin{Shaded}
\begin{Highlighting}[]
\NormalTok{response }\OperatorTok{=} \ControlFlowTok{await} \FunctionTok{fetch}\NormalTok{(}\KeywordTok{this}\OperatorTok{.}\AttributeTok{API}\OperatorTok{,}\NormalTok{ \{}
    \DataTypeTok{method}\OperatorTok{:} \StringTok{\textquotesingle{}POST\textquotesingle{}}\OperatorTok{,}
    \OperatorTok{...}
\NormalTok{\})}\OperatorTok{;}
\end{Highlighting}
\end{Shaded}

The URL stays the same. Only the \texttt{method} changes: POST to
create, PUT to update, DELETE to remove.

This is the key insight of the chapter: one URL, different
methods. Everything else follows from here. You act on it by
choosing the right HTTP method. The URL says \emph{what}, the
method says \emph{what to do with it}. Instead of inventing
separate URLs like \texttt{/api/publishers/create} and
\texttt{/api/publishers/delete}, we use a single URL and let the
HTTP method express the intent.

The separation between \texttt{load} and \texttt{render} is the
same in every class: \texttt{load} talks to the server,
\texttt{render} talks to the DOM. Once you understand
\texttt{Publisher}, you understand them all. \texttt{Book} adds
complexity (multiple authors, more fields), but the structure is
the same.

\newpage

\hypertarget{crud}{%
\chapter{CRUD}\label{crud}}

\hypertarget{the-sign-of-the-four}{%
\section{The sign of the four}\label{the-sign-of-the-four}}

CRUD stands for Create, Read, Update, Delete. Every data-driven
application does these four things. Nothing more, nothing less.

In the previous chapter we learned that HTTP has a method for
each operation. Here is the complete mapping, from the browser to
the database:

\begin{longtable}[]{@{}
  >{\raggedright\arraybackslash}p{\dimexpr (\columnwidth - 6\tabcolsep) * 0.1803 \relax}
  >{\raggedright\arraybackslash}p{\dimexpr (\columnwidth - 6\tabcolsep) * 0.2131 \relax}
  >{\raggedright\arraybackslash}p{\dimexpr (\columnwidth - 6\tabcolsep) * 0.1803 \relax}
  >{\raggedright\arraybackslash}p{\dimexpr (\columnwidth - 6\tabcolsep) * 0.4262 \relax}@{}}
\toprule
\begin{minipage}[b]{\linewidth}\raggedright
Operation
\end{minipage} & \begin{minipage}[b]{\linewidth}\raggedright
HTTP Method
\end{minipage} & \begin{minipage}[b]{\linewidth}\raggedright
SQL
\end{minipage} & \begin{minipage}[b]{\linewidth}\raggedright
Example
\end{minipage} \\
\midrule
\endhead
Create & \texttt{POST} & \texttt{INSERT} & add a new publisher \\
Read & \texttt{GET} & \texttt{SELECT} & list all publishers \\
Update & \texttt{PUT} & \texttt{UPDATE} & change a publisher's
name \\
Delete & \texttt{DELETE} & \texttt{DELETE} & remove a
publisher \\
\bottomrule
\end{longtable}

Three layers, one intention. The HTTP method arrives at the API
file (e.g.~\texttt{api/publishers.php}). In MVC architecture,
this file is called the \textbf{controller}: it receives the
request and decides what to do. It does not talk to the database
directly. Instead, it calls the \textbf{model}: a PHP class
(file) that lives in \texttt{src/} (outside \texttt{public/},
because it handles data, not HTTP requests). Each model ---
\texttt{Publisher.php}, \texttt{Category.php},
\texttt{Author.php}, \texttt{Book.php} --- implements the CRUD
methods internally: \texttt{insert}, \texttt{getAll},
\texttt{getById}, \texttt{update}, \texttt{delete}. None of them
talks to the database directly either. They all use
\texttt{DBMS.php}, a class we wrote that wraps PDO and exposes
simple methods: \texttt{fetchOne}, \texttt{fetchAll},
\texttt{execute}. Every model receives a \texttt{DBMS} instance
in its constructor and uses it for every query. One class for
database access, shared by all models. The controller calls the
model, the model calls DBMS, DBMS talks to SQLite, and the
controller returns the result as JSON.

\hypertarget{soft-delete-vs-hard-delete}{%
\section{Soft delete vs hard
delete}\label{soft-delete-vs-hard-delete}}

Bibliotheca uses a \texttt{status} column (1 = active, 0 =
disabled). When you disable a record, the row stays in the
database --- it just appears greyed out in the list. This is a
\textbf{soft delete}: a PUT that sets \texttt{status} to 0.

When you delete a record, it is gone. \texttt{DELETE\ FROM}
removes the row entirely. This is a \textbf{hard delete}:
permanent, irreversible.

Why both? In a real application, you often want to keep the data.
An author who is temporarily unavailable should not disappear
from the database --- their books still reference them. Disabling
is safer. Deleting is for cleanup.

\hypertarget{from-form-to-database}{%
\section{From form to database}\label{from-form-to-database}}

In a classic website, clicking a button loads a new page. In
Bibliotheca, JavaScript intercepts the action, calls
\texttt{fetch} with the right HTTP method, and updates the page
without reloading. Let us follow each operation step by step.

\hypertarget{create-post}{%
\subsection{Create, POST}\label{create-post}}

\begin{enumerate}
\def\labelenumi{\arabic{enumi}.}
\tightlist
\item
  The user fills the form and clicks Save.
\item
  JavaScript builds the \texttt{payload} object from the form
  fields.
\item
  \texttt{fetch} sends a POST with the payload as JSON body.
\item
  The controller validates and calls the model's \texttt{insert}.
\item
  The model executes \texttt{INSERT\ INTO} via PDO.
\item
  The controller returns the new record as JSON.
\item
  JavaScript redirects to the list page.
\end{enumerate}

\hypertarget{read-get}{%
\subsection{Read, GET}\label{read-get}}

\begin{enumerate}
\def\labelenumi{\arabic{enumi}.}
\tightlist
\item
  The page loads and JavaScript calls \texttt{fetch} with GET.
\item
  The controller calls the model's \texttt{getAll} or
  \texttt{getById}.
\item
  The model queries via PDO and returns an array.
\item
  The controller encodes the array as JSON.
\item
  JavaScript receives the data and calls \texttt{render} to build
  the table rows or populate the form fields.
\end{enumerate}

\hypertarget{update-put}{%
\subsection{Update, PUT}\label{update-put}}

\begin{enumerate}
\def\labelenumi{\arabic{enumi}.}
\tightlist
\item
  The user clicks Edit --- the browser goes to the form page with
  \texttt{?id=} in the URL.
\item
  \texttt{checkEdit} detects the ID and calls \texttt{load(id)}.
\item
  \texttt{load} fetches the record via GET, then \texttt{render}
  populates the form.
\item
  The user modifies the fields and clicks Save.
\item
  JavaScript builds the \texttt{payload} and adds the record ID.
\item
  \texttt{fetch} sends a PUT with the payload as JSON body.
\item
  The controller validates and calls the model's \texttt{update}.
\item
  JavaScript redirects to the list page.
\end{enumerate}

\hypertarget{delete-delete}{%
\subsection{Delete, DELETE}\label{delete-delete}}

\begin{enumerate}
\def\labelenumi{\arabic{enumi}.}
\tightlist
\item
  The user clicks Delete --- a confirmation dialog appears.
\item
  JavaScript builds the \texttt{payload} with the record ID.
\item
  \texttt{fetch} sends a DELETE with the payload as JSON body.
\item
  The controller calls the model's \texttt{delete}
  (\texttt{DELETE\ FROM} --- permanent).
\item
  JavaScript redirects to the list page.
\end{enumerate}

Notice how every operation follows the same shape: build the
payload, choose the method, call \texttt{fetch}, handle the
response. The only things that change are the HTTP method and the
data inside the payload.

\newpage

\hypertarget{validation}{%
\chapter{Validation}\label{validation}}

\hypertarget{three-lines-of-defense}{%
\section{Three lines of defense}\label{three-lines-of-defense}}

Validation is not optional. It is part of the craft.

\begin{enumerate}
\def\labelenumi{\arabic{enumi}.}
\item
  \textbf{HTML} --- we set constraints directly on the form
  fields: \texttt{required}, \texttt{maxlength},
  \texttt{type="number"}, \texttt{min}, \texttt{max}. The browser
  enforces these before any code runs. We add \texttt{novalidate}
  to the \texttt{\textless{}form\textgreater{}} tag so the
  browser does not show its own popups --- we want to control the
  error messages ourselves.
\item
  \textbf{JavaScript} --- we delegate the interactive checks to
  the \texttt{validate()} method, which runs before sending the
  request. Inline error messages appear under each field, red
  borders highlight invalid inputs. This is for the user
  experience: fast feedback, no round-trip to the server.
\item
  \textbf{PHP} --- we delegate the final check to the server.
  Never trust the client. The server validates everything again,
  because anyone can bypass JavaScript with a single line in the
  browser console.
\end{enumerate}

Each layer validates as if it were the only one. HTML does not
know JavaScript exists. JavaScript does not assume the server
will catch what it missed. The server does not trust anything
that came before it. They each see the data from their own point
of view, and they each reject bad input independently. When you
write the validation for one layer, forget the others exist. Do
not think ``HTML already checks this'' or ``the server will catch
it anyway''. Each layer must stand on its own and do its job to
the best of its capabilities.

\hypertarget{what-we-validate-in-our-project}{%
\section{What we validate in our
project}\label{what-we-validate-in-our-project}}

\begin{itemize}
\tightlist
\item
  \textbf{Required fields} --- cannot be empty after trimming.
\item
  \textbf{String length} --- names capped at 100 characters,
  titles at 255.
\item
  \textbf{Numeric ranges} --- pages between 1 and 99999,
  published year from 1000 to the current year.
\item
  \textbf{Date format and range} --- author birthdate must be a
  valid \texttt{YYYY-MM-DD}, year \textgreater= 1000, not in the
  future.
\item
  \textbf{Duplicates} --- server-side check with the model's
  \texttt{exists()} method. Returns HTTP 409 (Conflict) if the
  record already exists.
\item
  \textbf{Dependencies} --- returns HTTP 422 (Unprocessable
  Entity) if the record has linked books and cannot be deleted or
  disabled.
\item
  \textbf{Input filtering} --- numeric fields only accept digit
  keys.
\end{itemize}

\hypertarget{input-normalization}{%
\section{Input normalization}\label{input-normalization}}

Before data reaches the database, the server normalizes it. This
is not cosmetic --- it is data entry control. If you let
``eINAUDI'', '' Einaudi ``, and''einaudi'' into the database, you
have three records for the same publisher. Same with numbers: if
a field expects an integer (like pages or published year), then
only digits should be allowed. Not letters, not decimals, not
special characters. The type of number matters: an integer is not
a float, a year is not a price. Look at the column type in the
database table: if it says \texttt{INTEGER}, the input field
should only accept digits. A winning strategy: define what you
expect based on where the data will be stored, and reject
everything else at the point of entry.

\begin{itemize}
\tightlist
\item
  \texttt{trim()} --- no leading or trailing spaces.
\item
  \texttt{ucwords(strtolower())} --- ``eINAUDI'' becomes
  ``Einaudi''.
\end{itemize}

The user types whatever they want. The server stores it in a
consistent format, every time.

\hypertarget{sanitization}{%
\section{Sanitization}\label{sanitization}}

Formatting makes input look right. Sanitization makes it
\emph{safe}.

Every string that enters the API goes through three filters
before anything else happens:

\begin{Shaded}
\begin{Highlighting}[]
\VariableTok{$name} \OperatorTok{=} \FunctionTok{mb\_substr}\NormalTok{(}\FunctionTok{strip\_tags}\NormalTok{(}\FunctionTok{trim}\NormalTok{(}\VariableTok{$data}\NormalTok{[}\StringTok{\textquotesingle{}name\textquotesingle{}}\NormalTok{] }\OperatorTok{??} \StringTok{\textquotesingle{}\textquotesingle{}}\NormalTok{))}\OtherTok{,} \DecValTok{0}\OtherTok{,} \DecValTok{100}\NormalTok{)}\OtherTok{;}
\end{Highlighting}
\end{Shaded}

\begin{itemize}
\tightlist
\item
  \textbf{\texttt{trim()}} --- removes leading and trailing
  whitespace.
\item
  \textbf{\texttt{strip\_tags()}} --- removes any HTML or PHP
  tags. A name like
  \texttt{\textless{}script\textgreater{}alert(\textquotesingle{}xss\textquotesingle{})\textless{}/script\textgreater{}}
  becomes
  \texttt{alert(\textquotesingle{}xss\textquotesingle{})}. The
  dangerous part is gone before the value reaches the database.
\item
  \textbf{\texttt{mb\_substr()}} --- enforces a maximum length.
  Even if someone sends a megabyte of text, only the first 100
  characters survive.
\end{itemize}

The order matters: trim first, strip second, truncate last. And
this happens on the server --- never trust the client's
\texttt{maxlength} attribute alone.

\hypertarget{error-display-how-do-we-tell-the-user-something-is-wrong}{%
\section{Error display: how do we tell the user something is
wrong?}\label{error-display-how-do-we-tell-the-user-something-is-wrong}}

No \texttt{alert()}. We use inline messages:

\begin{Shaded}
\begin{Highlighting}[]
\KeywordTok{\textless{}span} \ErrorTok{class}\OtherTok{=}\StringTok{"error"} \ErrorTok{id}\OtherTok{=}\StringTok{"publisher{-}name{-}error"}\KeywordTok{\textgreater{}\textless{}/span\textgreater{}}
\end{Highlighting}
\end{Shaded}

Each input field has a matching
\texttt{\textless{}span\textgreater{}} for its error message. The
JavaScript class has two methods for this:

\begin{Shaded}
\begin{Highlighting}[]
\FunctionTok{showError}\NormalTok{(fieldId}\OperatorTok{,}\NormalTok{ message) \{}
    \KeywordTok{const}\NormalTok{ input }\OperatorTok{=} \BuiltInTok{document}\OperatorTok{.}\FunctionTok{querySelector}\NormalTok{(}\StringTok{\textquotesingle{}\#\textquotesingle{}} \OperatorTok{+}\NormalTok{ fieldId)}\OperatorTok{;}
    \KeywordTok{const}\NormalTok{ error }\OperatorTok{=} \BuiltInTok{document}\OperatorTok{.}\FunctionTok{querySelector}\NormalTok{(}
        \StringTok{\textquotesingle{}\#\textquotesingle{}} \OperatorTok{+}\NormalTok{ fieldId }\OperatorTok{+} \StringTok{\textquotesingle{}{-}error\textquotesingle{}}
\NormalTok{    )}\OperatorTok{;}
\NormalTok{    input}\OperatorTok{.}\AttributeTok{classList}\OperatorTok{.}\FunctionTok{add}\NormalTok{(}\StringTok{\textquotesingle{}invalid\textquotesingle{}}\NormalTok{)}\OperatorTok{;}
\NormalTok{    error}\OperatorTok{.}\AttributeTok{textContent} \OperatorTok{=}\NormalTok{ message}\OperatorTok{;}
\NormalTok{\}}
\end{Highlighting}
\end{Shaded}

\begin{Shaded}
\begin{Highlighting}[]
\FunctionTok{clearErrors}\NormalTok{() \{}
    \KeywordTok{const}\NormalTok{ errors }\OperatorTok{=} \BuiltInTok{document}\OperatorTok{.}\FunctionTok{querySelectorAll}\NormalTok{(}\StringTok{\textquotesingle{}.error\textquotesingle{}}\NormalTok{)}\OperatorTok{;}
    \ControlFlowTok{for}\NormalTok{ (}\KeywordTok{let}\NormalTok{ i }\OperatorTok{=} \DecValTok{0}\OperatorTok{;}\NormalTok{ i }\OperatorTok{\textless{}}\NormalTok{ errors}\OperatorTok{.}\AttributeTok{length}\OperatorTok{;}\NormalTok{ i}\OperatorTok{++}\NormalTok{) \{}
\NormalTok{        errors[i]}\OperatorTok{.}\AttributeTok{textContent} \OperatorTok{=} \StringTok{\textquotesingle{}\textquotesingle{}}\OperatorTok{;}
\NormalTok{    \}}
    \KeywordTok{const}\NormalTok{ invalids }\OperatorTok{=} \BuiltInTok{document}\OperatorTok{.}\FunctionTok{querySelectorAll}\NormalTok{(}\StringTok{\textquotesingle{}.invalid\textquotesingle{}}\NormalTok{)}\OperatorTok{;}
    \ControlFlowTok{for}\NormalTok{ (}\KeywordTok{let}\NormalTok{ i }\OperatorTok{=} \DecValTok{0}\OperatorTok{;}\NormalTok{ i }\OperatorTok{\textless{}}\NormalTok{ invalids}\OperatorTok{.}\AttributeTok{length}\OperatorTok{;}\NormalTok{ i}\OperatorTok{++}\NormalTok{) \{}
\NormalTok{        invalids[i]}\OperatorTok{.}\AttributeTok{classList}\OperatorTok{.}\FunctionTok{remove}\NormalTok{(}\StringTok{\textquotesingle{}invalid\textquotesingle{}}\NormalTok{)}\OperatorTok{;}
\NormalTok{    \}}
\NormalTok{\}}
\end{Highlighting}
\end{Shaded}

\texttt{showError} marks one field as invalid and shows the
message. \texttt{clearErrors} resets all fields before a new
validation pass. The \texttt{validate()} method calls
\texttt{clearErrors} first, then checks each field and calls
\texttt{showError} for every problem it finds.

\hypertarget{server-errors}{%
\section{Server errors}\label{server-errors}}

When the server rejects the input (409 for duplicates, 422 for
dependencies, 400 for bad data), JavaScript reads the response
and shows the message inline --- same place, same pattern:

\begin{Shaded}
\begin{Highlighting}[]
\ControlFlowTok{if}\NormalTok{ (}\OperatorTok{!}\NormalTok{response}\OperatorTok{.}\AttributeTok{ok}\NormalTok{) \{}
    \KeywordTok{const}\NormalTok{ result }\OperatorTok{=} \ControlFlowTok{await}\NormalTok{ response}\OperatorTok{.}\FunctionTok{json}\NormalTok{()}\OperatorTok{;}
    \KeywordTok{this}\OperatorTok{.}\FunctionTok{showError}\NormalTok{(}\StringTok{\textquotesingle{}publisher{-}name\textquotesingle{}}\OperatorTok{,}\NormalTok{ result}\OperatorTok{.}\AttributeTok{error}\NormalTok{)}\OperatorTok{;}
    \ControlFlowTok{return}\OperatorTok{;}
\NormalTok{\}}
\end{Highlighting}
\end{Shaded}

Whether the error comes from client-side validation or from the
server, the user sees it in the same way.

\hypertarget{exceptions-who-catches-what}{%
\section{Exceptions: who catches
what}\label{exceptions-who-catches-what}}

Validation handles \emph{expected} problems --- empty fields,
duplicates. But what about \emph{unexpected} failures? The
database file is corrupted. The disk is full. A table was
dropped. These are exceptions.

In PHP, PDO is configured with \texttt{ERRMODE\_EXCEPTION}: when
something goes wrong, it throws a \texttt{PDOException}. The
question is: who catches it?

\textbf{Not DBMS.} The database wrapper does not know \emph{how}
to handle the error. Should it return null? An empty array? Log
something? It depends on who is calling. A web API needs a JSON
response. A CLI script needs a console message. A test needs an
assertion failure. DBMS cannot know any of this, so it lets the
exception rise.

\textbf{Not the Model.} Publisher, Book, Author --- they have the
same problem. They do not know the context. They pass the
exception upward.

\textbf{The controller catches it.} The API file is the only
place that knows it must respond with JSON and an HTTP status
code:

\begin{Shaded}
\begin{Highlighting}[]
\ControlFlowTok{try}\NormalTok{ \{}
    \VariableTok{$db} \OperatorTok{=} \KeywordTok{new} \ConstantTok{DBMS}\NormalTok{(}
        \ConstantTok{\_\_DIR\_\_} \OperatorTok{.} \StringTok{\textquotesingle{}/../../sql/bibliotheca.db\textquotesingle{}}
\NormalTok{    )}\OtherTok{;}
    \VariableTok{$publisher} \OperatorTok{=} \KeywordTok{new} \ConstantTok{P}\NormalTok{ublisher(}\VariableTok{$db}\NormalTok{)}\OtherTok{;}
    \CommentTok{// ... all the logic ...}
\NormalTok{\} }\ControlFlowTok{catch}\NormalTok{ (\textbackslash{}}\BuiltInTok{Exception} \VariableTok{$e}\NormalTok{) \{}
    \FunctionTok{http\_response\_code}\NormalTok{(}\DecValTok{500}\NormalTok{)}\OtherTok{;}
    \KeywordTok{echo} \FunctionTok{json\_encode}\NormalTok{([}\StringTok{\textquotesingle{}error\textquotesingle{}}\NormalTok{ =\textgreater{} }\StringTok{\textquotesingle{}Internal server error\textquotesingle{}}\NormalTok{])}\OtherTok{;}
\NormalTok{\}}
\end{Highlighting}
\end{Shaded}

\textbf{JavaScript catches too.} Network failures (server down,
timeout) throw exceptions at the \texttt{fetch} level:

\begin{Shaded}
\begin{Highlighting}[]
\ControlFlowTok{try}\NormalTok{ \{}
    \KeywordTok{const}\NormalTok{ response }\OperatorTok{=} \ControlFlowTok{await} \FunctionTok{fetch}\NormalTok{(}\KeywordTok{this}\OperatorTok{.}\AttributeTok{API}\NormalTok{)}\OperatorTok{;}
    \ControlFlowTok{if}\NormalTok{ (}\OperatorTok{!}\NormalTok{response}\OperatorTok{.}\AttributeTok{ok}\NormalTok{) \{}
        \ControlFlowTok{throw} \KeywordTok{new} \BuiltInTok{Error}\NormalTok{(response}\OperatorTok{.}\AttributeTok{statusText}\NormalTok{)}\OperatorTok{;}
\NormalTok{    \}}
    \KeywordTok{const}\NormalTok{ data }\OperatorTok{=} \ControlFlowTok{await}\NormalTok{ response}\OperatorTok{.}\FunctionTok{json}\NormalTok{()}\OperatorTok{;}
    \KeywordTok{this}\OperatorTok{.}\FunctionTok{render}\NormalTok{(data)}\OperatorTok{;}
\NormalTok{\} }\ControlFlowTok{catch}\NormalTok{ (error) \{}
    \FunctionTok{alert}\NormalTok{(}\StringTok{\textquotesingle{}Unable to load data.\textquotesingle{}}\NormalTok{)}\OperatorTok{;}
\NormalTok{\}}
\end{Highlighting}
\end{Shaded}

The principle: \textbf{catch where you can act.} The lower layers
report the problem (by throwing). The upper layers handle it (by
catching). An exception is a signal that travels upward --- like
pain traveling from an organ to the brain. The kidney does not
decide how to inform the patient. The brain does.

\newpage

\hypertarget{permissions}{%
\chapter{Permissions}\label{permissions}}

\hypertarget{the-problem}{%
\section{The problem}\label{the-problem}}

Your code works from the terminal. You open the browser and get a
blank page, or worse, a 500 Internal Server Error. The reason is
almost always permissions.

Apache runs as a user called \texttt{www-data}. Your files belong
to your user (e.g., \texttt{bc}). If \texttt{www-data} cannot
read a file, the page is blank. If it cannot write, the database
operations fail silently.

\hypertarget{how-linux-permissions-work}{%
\section{How Linux permissions
work}\label{how-linux-permissions-work}}

Every file has three sets of permissions: owner, group, others.

\begin{verbatim}
-rw-r--r-- 1 youruser youruser 28672 bibliotheca.db
\end{verbatim}

\begin{itemize}
\tightlist
\item
  \texttt{rw-} --- owner (\texttt{youruser}) can read and write.
\item
  \texttt{r-\/-} --- group (\texttt{youruser}) can only read.
\item
  \texttt{r-\/-} --- others can only read.
\end{itemize}

Apache (\texttt{www-data}) is ``others'' here --- it can read but
not write. That is why SELECT works but INSERT fails.

\hypertarget{the-fix}{%
\section{The fix}\label{the-fix}}

Give the group \texttt{www-data} write access to the database
file and its directory (SQLite needs the directory for its
journal file):

\begin{Shaded}
\begin{Highlighting}[]
\FunctionTok{sudo}\NormalTok{ chgrp www{-}data /var/www/html/bibliotheca/sql/bibliotheca.db}
\FunctionTok{sudo}\NormalTok{ chmod 664 /var/www/html/bibliotheca/sql/bibliotheca.db}
\FunctionTok{sudo}\NormalTok{ chgrp www{-}data /var/www/html/bibliotheca/sql/}
\FunctionTok{sudo}\NormalTok{ chmod 775 /var/www/html/bibliotheca/sql/}
\end{Highlighting}
\end{Shaded}

\begin{itemize}
\tightlist
\item
  \texttt{664} --- owner reads/writes, group reads/writes, others
  read.
\item
  \texttt{775} --- same for directories, plus execute (needed to
  list contents).
\end{itemize}

Check the result with \texttt{ls\ -la}:

\begin{Shaded}
\begin{Highlighting}[]
\FunctionTok{ls} \AttributeTok{{-}la}\NormalTok{ sql/}
\end{Highlighting}
\end{Shaded}

The terminal is your best friend for diagnosing permission
problems. Always verify with \texttt{ls\ -la} after changing
permissions. Get comfortable with the command line --- it should
always be your first choice. A production server has no graphical
interface. There is no file manager, no right-click, no
properties dialog. There is only the terminal. The sooner you
make it your natural environment, the better.

\hypertarget{every-time-you-recreate-the-database}{%
\section{Every time you recreate the
database}\label{every-time-you-recreate-the-database}}

When you delete and recreate \texttt{bibliotheca.db}, the new
file belongs to your user with default permissions. You must set
them again:

\begin{Shaded}
\begin{Highlighting}[]
\FunctionTok{rm}\NormalTok{ bibliotheca.db}
\ExtensionTok{sqlite3}\NormalTok{ bibliotheca.db }\StringTok{".read ddl/create\_table.sql"}
\ExtensionTok{sqlite3}\NormalTok{ bibliotheca.db }\StringTok{".read dml/insert.sql"}
\FunctionTok{sudo}\NormalTok{ chgrp www{-}data bibliotheca.db}
\FunctionTok{sudo}\NormalTok{ chmod 664 bibliotheca.db}
\end{Highlighting}
\end{Shaded}

This is easy to forget. And when you forget, the application
reads data fine but fails on any write operation with:

\begin{verbatim}
SQLSTATE[HY000]: General error: 8 attempt to write a readonly database
\end{verbatim}

If you see this error, check the permissions first. Nine times
out of ten, that is the problem.

\hypertarget{why-public-matters}{%
\section{\texorpdfstring{Why \texttt{public/}
matters}{Why public/ matters}}\label{why-public-matters}}

Notice the project structure: the \texttt{src/} folder (models,
DBMS) and the \texttt{sql/} folder (database) are outside
\texttt{public/}. Only the files inside \texttt{public/} are
reachable from the browser.

This is a permission boundary too. Even if someone guesses the
path to \texttt{sql/bibliotheca.db}, Apache will not serve it ---
the \texttt{.htaccess} rewrite rule sends everything through
\texttt{index.php}, and the router only accepts whitelisted
routes.

The file system layout is your first line of defense. Keep
sensitive files outside the web root.

\newpage

\hypertarget{debugging}{%
\chapter{Debugging}\label{debugging}}

\hypertarget{the-blank-page}{%
\section{The blank page}\label{the-blank-page}}

The most frustrating thing in web development: your code looks
correct, but the browser shows nothing. No error, no message,
just white. This means PHP crashed before producing any output.

The browser cannot help you here. You need to look elsewhere.

\hypertarget{think-in-layers}{%
\section{Think in layers}\label{think-in-layers}}

When something breaks, remember the architecture from Chapter 06:
JavaScript calls the controller, the controller calls the model,
the model calls DBMS, DBMS talks to SQLite. The bug lives in one
of these layers. Your job is to find which one.

Start from the outside and work inward:

\begin{enumerate}
\def\labelenumi{\arabic{enumi}.}
\tightlist
\item
  Is the browser showing a JavaScript error? → \textbf{Console}
\item
  Is the server returning an error code? → \textbf{Network tab}
\item
  Is PHP crashing? → \textbf{Apache error log}
\item
  Is the SQL wrong? → \textbf{SQLite command line}
\end{enumerate}

Be methodical. Narrow it down.

\hypertarget{apache-error-log}{%
\section{Apache error log}\label{apache-error-log}}

The server-side truth lives here:

\begin{Shaded}
\begin{Highlighting}[]
\FunctionTok{sudo}\NormalTok{ tail }\AttributeTok{{-}20}\NormalTok{ /var/log/apache2/error.log}
\end{Highlighting}
\end{Shaded}

This log captures every PHP fatal error, every uncaught
exception, every permission problem. When the browser shows a
500, this log tells you why.

For real-time monitoring, use \texttt{-f} --- the log updates as
requests come in:

\begin{Shaded}
\begin{Highlighting}[]
\FunctionTok{sudo}\NormalTok{ tail }\AttributeTok{{-}f}\NormalTok{ /var/log/apache2/error.log}
\end{Highlighting}
\end{Shaded}

Keep this running in a terminal while you test in the browser.
You will see errors the moment they happen.

Common errors you will find:

\begin{itemize}
\tightlist
\item
  \texttt{could\ not\ find\ driver} --- a PHP extension is
  missing.
\item
  \texttt{attempt\ to\ write\ a\ readonly\ database} ---
  permissions (Chapter 08).
\item
  \texttt{Uncaught\ PDOException} --- a SQL error, with the exact
  line number.
\end{itemize}

\hypertarget{curl-testing-the-api}{%
\section{curl: testing the API}\label{curl-testing-the-api}}

Before blaming the frontend, test the backend directly:

\begin{Shaded}
\begin{Highlighting}[]
\ExtensionTok{curl}\NormalTok{ http://localhost/bibliotheca/public/api/publishers.php}
\end{Highlighting}
\end{Shaded}

If this returns JSON, the backend works. If it returns nothing or
an error, the problem is server-side.

For POST requests, remember that the controller requires a CSRF
token (Chapter 10). Without it, you will get a 403. For a quick
diagnostic, you can test a GET first to confirm the API is alive,
then check the Network tab for POST issues.

Use \texttt{-v} for verbose output --- it shows HTTP status
codes, headers, and the full request/response cycle:

\begin{Shaded}
\begin{Highlighting}[]
\ExtensionTok{curl} \AttributeTok{{-}v}\NormalTok{ http://localhost/bibliotheca/public/api/publishers.php}
\end{Highlighting}
\end{Shaded}

The \texttt{-v} flag is one of the most useful debugging tools
you have. It shows you exactly what the browser and server are
saying to each other.

\hypertarget{browser-developer-tools}{%
\section{Browser developer tools}\label{browser-developer-tools}}

Press F12. Three tabs matter:

\begin{itemize}
\tightlist
\item
  \textbf{Console} --- JavaScript errors appear here. If
  \texttt{fetch} fails or your code has a syntax error, this is
  where you see it.
\item
  \textbf{Network} --- every HTTP request the browser makes.
  Click on a request to see the status code, request body, and
  response body. If a fetch returns 500, click on it to see the
  server's error message.
\item
  \textbf{Elements} --- the live DOM. Useful to check if
  JavaScript actually built the HTML you expected. Remember the
  DOM diagram from Chapter 05: what you see here is the tree, not
  the source file.
\end{itemize}

\hypertarget{php-debugging-tools}{%
\section{PHP debugging tools}\label{php-debugging-tools}}

When you need to inspect a value inside your PHP code, two tools:

\textbf{\texttt{error\_log()}} --- writes to the Apache error log
without breaking the response:

\begin{Shaded}
\begin{Highlighting}[]
\FunctionTok{error\_log}\NormalTok{(}\StringTok{\textquotesingle{}Publisher name: \textquotesingle{}} \OperatorTok{.} \VariableTok{$name}\NormalTok{)}\OtherTok{;}
\FunctionTok{error\_log}\NormalTok{(}\FunctionTok{print\_r}\NormalTok{(}\VariableTok{$data}\OtherTok{,} \KeywordTok{true}\NormalTok{))}\OtherTok{;}
\end{Highlighting}
\end{Shaded}

Check the output with
\texttt{tail\ -f\ /var/log/apache2/error.log}. This is safe to
use in API controllers --- it does not corrupt the JSON response.

\textbf{\texttt{var\_dump()}} --- prints the value directly to
the output:

\begin{Shaded}
\begin{Highlighting}[]
\FunctionTok{var\_dump}\NormalTok{(}\VariableTok{$data}\NormalTok{)}\OtherTok{;}
\ControlFlowTok{exit}\OtherTok{;}
\end{Highlighting}
\end{Shaded}

Useful for quick inspection, but it breaks the JSON response. Use
it only for temporary debugging, never in production code. Remove
it when you are done.

\hypertarget{php-from-the-command-line}{%
\section{PHP from the command
line}\label{php-from-the-command-line}}

You can run a PHP file directly to see errors without Apache:

\begin{Shaded}
\begin{Highlighting}[]
\ExtensionTok{php}\NormalTok{ /var/www/html/bibliotheca/public/api/publishers.php}
\end{Highlighting}
\end{Shaded}

This bypasses Apache and permissions. If it works here but not in
the browser, the problem is Apache configuration or permissions
(Chapter 08).

\hypertarget{sqlite-from-the-command-line}{%
\section{SQLite from the command
line}\label{sqlite-from-the-command-line}}

You can inspect the database directly:

\begin{Shaded}
\begin{Highlighting}[]
\ExtensionTok{sqlite3}\NormalTok{ sql/bibliotheca.db}
\end{Highlighting}
\end{Shaded}

Useful commands inside the SQLite shell:

\begin{Shaded}
\begin{Highlighting}[]
\CommentTok{{-}{-} List all tables}
\NormalTok{.}\KeywordTok{tables}

\CommentTok{{-}{-} Show the schema of a table}
\NormalTok{PRAGMA table\_info(book);}

\CommentTok{{-}{-} Show all indexes in the database}
\NormalTok{.indices}

\CommentTok{{-}{-} Show the indexes of a specific table}
\NormalTok{PRAGMA index\_list(book);}

\CommentTok{{-}{-} Show which columns an index covers}
\NormalTok{PRAGMA index\_info(idx\_book\_publisher);}

\CommentTok{{-}{-} Check if foreign keys are enforced}
\NormalTok{PRAGMA foreign\_keys;}

\CommentTok{{-}{-} Show the CREATE statement for a table}
\NormalTok{.}\KeywordTok{schema}\NormalTok{ book}

\CommentTok{{-}{-} Run a quick query}
\KeywordTok{SELECT} \OperatorTok{*} \KeywordTok{FROM}\NormalTok{ publisher }\KeywordTok{ORDER} \KeywordTok{BY}\NormalTok{ name;}
\end{Highlighting}
\end{Shaded}

\texttt{PRAGMA\ table\_info} is especially useful: it shows the
column name, type, whether it allows NULL, its default value, and
whether it is part of the primary key. Learn to use it --- it is
faster than reading the schema file.

\hypertarget{the-golden-rule}{%
\section{The golden rule}\label{the-golden-rule}}

When something breaks, resist the urge to change random things
hoping the problem goes away. Read the error message. Follow the
chain: browser console → network tab → Apache log → PHP command
line → SQLite shell. The answer is always in one of these places.
Your job is not to fix --- it is to \emph{find}. Once you find
it, the fix is usually obvious.

\newpage

\hypertarget{security}{%
\chapter{Security}\label{security}}

\hypertarget{what-we-already-have}{%
\section{What we already have}\label{what-we-already-have}}

Previous chapters covered the essentials:

\begin{itemize}
\tightlist
\item
  \textbf{SQL injection} --- prepared statements with named
  parameters (Chapter 04).
\item
  \textbf{XSS} --- \texttt{htmlspecialchars()} in PHP,
  \texttt{textContent} in JavaScript (Chapter 07).
\item
  \textbf{Input validation} --- server-side checks on every field
  (Chapter 07).
\item
  \textbf{File permissions} --- keeping sensitive files outside
  the web root (Chapter 08).
\end{itemize}

These are non-negotiable. They are part of the code you have
already written. This chapter adds the layers that protect
\emph{around} your code.

\hypertarget{csrf-protecting-write-operations}{%
\section{CSRF --- protecting write
operations}\label{csrf-protecting-write-operations}}

In Chapter 05 we saw that JavaScript sends an
\texttt{X-CSRF-Token} header with every POST, PUT, and DELETE
request. But what is CSRF, and why do we need that token?

\hypertarget{the-attack}{%
\subsection{The attack}\label{the-attack}}

Imagine you have Bibliotheca open in your browser. A malicious
website in another tab contains hidden JavaScript that sends a
POST request to your \texttt{api/publishers.php} with
\texttt{\{"name":\ "HACKED"\}}. The browser sends it --- from its
perspective, it is just another HTTP request to localhost, and it
attaches your session cookie automatically.

This is \textbf{CSRF} (Cross-Site Request Forgery): a malicious
site tricks the browser into making requests to \emph{your}
application, using \emph{your} session.

\hypertarget{the-defense-a-secret-token}{%
\subsection{The defense: a secret
token}\label{the-defense-a-secret-token}}

The solution is a \textbf{CSRF token}: a long random string that
only our pages know. The malicious site cannot read it because
the browser's \textbf{Same-Origin Policy} prevents one site from
reading another site's content.

The flow:

\begin{enumerate}
\def\labelenumi{\arabic{enumi}.}
\item
  When PHP renders the page (\texttt{index.php}), it generates a
  random token and stores it in the session:

\begin{Shaded}
\begin{Highlighting}[]
\VariableTok{$\_SESSION}\NormalTok{[}\StringTok{\textquotesingle{}csrf\_token\textquotesingle{}}\NormalTok{] }\OperatorTok{=} \FunctionTok{bin2hex}\NormalTok{(}\FunctionTok{random\_bytes}\NormalTok{(}\DecValTok{32}\NormalTok{))}\OtherTok{;}
\end{Highlighting}
\end{Shaded}
\item
  The token is injected into the HTML as a
  \texttt{\textless{}meta\textgreater{}} tag:

\begin{Shaded}
\begin{Highlighting}[]
\KeywordTok{\textless{}meta} \ErrorTok{name}\OtherTok{=}\StringTok{"csrf{-}token"} \ErrorTok{content}\OtherTok{=}\StringTok{"a7f3b9c2e1d4..."}\KeywordTok{\textgreater{}}
\end{Highlighting}
\end{Shaded}
\item
  JavaScript reads the token and includes it in every POST, PUT,
  and DELETE request --- the same header we saw in Chapter 05:

\begin{Shaded}
\begin{Highlighting}[]
\NormalTok{headers}\OperatorTok{:}\NormalTok{ \{}
    \StringTok{\textquotesingle{}Content{-}Type\textquotesingle{}}\OperatorTok{:} \StringTok{\textquotesingle{}application/json\textquotesingle{}}\OperatorTok{,}
    \StringTok{\textquotesingle{}X{-}CSRF{-}Token\textquotesingle{}}\OperatorTok{:} \BuiltInTok{document}\OperatorTok{.}\FunctionTok{querySelector}\NormalTok{(}
        \StringTok{\textquotesingle{}meta[name="csrf{-}token"]\textquotesingle{}}
\NormalTok{    )}\OperatorTok{.}\AttributeTok{content}
\NormalTok{\}}
\end{Highlighting}
\end{Shaded}
\item
  The controller (the API file) verifies the token before
  processing the request:

\begin{Shaded}
\begin{Highlighting}[]
\ConstantTok{C}\NormalTok{srf::verify()}\OtherTok{;}  \CommentTok{// 403 if token is missing or wrong}
\end{Highlighting}
\end{Shaded}
\end{enumerate}

The malicious site cannot read the
\texttt{\textless{}meta\textgreater{}} tag from our page
(Same-Origin Policy), so it cannot send the correct token. The
controller rejects its request with \texttt{403\ Forbidden}.

\hypertarget{the-csrf-class}{%
\subsection{The Csrf class}\label{the-csrf-class}}

The \texttt{Csrf} class (\texttt{src/Csrf.php}) has three
methods:

\begin{itemize}
\tightlist
\item
  \textbf{\texttt{Csrf::start()}} --- starts a PHP session (with
  secure cookie settings).
\item
  \textbf{\texttt{Csrf::token()}} --- returns the token,
  generating one if needed.
\item
  \textbf{\texttt{Csrf::verify()}} --- compares the
  \texttt{X-CSRF-Token} header against the session; exits with
  403 on mismatch.
\end{itemize}

GET requests do not need CSRF protection --- they only read data,
they never modify it. Only POST, PUT, and DELETE are checked.

\hypertarget{verify-it-works-from-the-terminal}{%
\subsection{Verify it works from the
terminal}\label{verify-it-works-from-the-terminal}}

\begin{Shaded}
\begin{Highlighting}[]
\CommentTok{\# Without token: 403}
\ExtensionTok{curl} \AttributeTok{{-}s} \AttributeTok{{-}o}\NormalTok{ /dev/null }\AttributeTok{{-}w} \StringTok{"\%\{http\_code\}"} \DataTypeTok{\textbackslash{}}
     \AttributeTok{{-}X}\NormalTok{ POST }\DataTypeTok{\textbackslash{}}
\NormalTok{     http://localhost/bibliotheca/public/api/publishers.php }\DataTypeTok{\textbackslash{}}
     \AttributeTok{{-}H} \StringTok{"Content{-}Type: application/json"} \DataTypeTok{\textbackslash{}}
     \AttributeTok{{-}d} \StringTok{\textquotesingle{}\{"name": "Test"\}\textquotesingle{}}
\end{Highlighting}
\end{Shaded}

\begin{Shaded}
\begin{Highlighting}[]
\CommentTok{\# With token: 200}
\VariableTok{TOKEN}\OperatorTok{=}\VariableTok{$(}\ExtensionTok{curl} \AttributeTok{{-}s} \AttributeTok{{-}c}\NormalTok{ /tmp/c.txt }\DataTypeTok{\textbackslash{}}
\NormalTok{    http://localhost/bibliotheca/public/ }\DataTypeTok{\textbackslash{}}
    \KeywordTok{|} \FunctionTok{grep}\NormalTok{ csrf{-}token }\DataTypeTok{\textbackslash{}}
    \KeywordTok{|} \FunctionTok{sed} \StringTok{\textquotesingle{}s/.*content="\textbackslash{}([\^{}"]*\textbackslash{})".*/\textbackslash{}1/\textquotesingle{}}\VariableTok{)}

\ExtensionTok{curl} \AttributeTok{{-}s} \AttributeTok{{-}o}\NormalTok{ /dev/null }\AttributeTok{{-}w} \StringTok{"\%\{http\_code\}"} \DataTypeTok{\textbackslash{}}
     \AttributeTok{{-}b}\NormalTok{ /tmp/c.txt }\AttributeTok{{-}X}\NormalTok{ POST }\DataTypeTok{\textbackslash{}}
\NormalTok{     http://localhost/bibliotheca/public/api/publishers.php }\DataTypeTok{\textbackslash{}}
     \AttributeTok{{-}H} \StringTok{"Content{-}Type: application/json"} \DataTypeTok{\textbackslash{}}
     \AttributeTok{{-}H} \StringTok{"X{-}CSRF{-}Token: }\VariableTok{$TOKEN}\StringTok{"} \DataTypeTok{\textbackslash{}}
     \AttributeTok{{-}d} \StringTok{\textquotesingle{}\{"name": "Test"\}\textquotesingle{}}
\end{Highlighting}
\end{Shaded}

The first request has no token --- the server rejects it. The
second extracts the token from the page and sends it --- the
server accepts it. Try both from your terminal.

\hypertarget{sessions-and-session-fixation}{%
\section{Sessions and session
fixation}\label{sessions-and-session-fixation}}

HTTP is stateless: each request is independent, the server does
not remember who you are. \textbf{Sessions} solve this. When you
visit the application, PHP creates a session --- a small file on
the server identified by a random ID. The browser stores this ID
in a cookie and sends it with every request. That is how the
server knows you are still you.

When you log in, PHP writes your user ID into the session. From
that point, every request carries your identity.

A \textbf{session fixation} attack exploits this: the attacker
gives you a session ID they already know (via a crafted URL or a
cookie), then waits for you to log in. Now they share your
authenticated session.

The defense is simple: regenerate the session ID after any
privilege change.

\begin{Shaded}
\begin{Highlighting}[]
\FunctionTok{session\_start}\NormalTok{()}\OtherTok{;}
\FunctionTok{session\_regenerate\_id}\NormalTok{(}\KeywordTok{true}\NormalTok{)}\OtherTok{;}  \CommentTok{// true = delete the old session file}
\end{Highlighting}
\end{Shaded}

The \texttt{true} parameter is important. Without it, the old
session file survives and the attacker's ID still works.

Bibliotheca calls \texttt{session\_regenerate\_id(true)} in
\texttt{Auth::login()} immediately after verifying credentials.
Our \texttt{Csrf::start()} calls \texttt{session\_start()} with
secure cookie settings:

\begin{Shaded}
\begin{Highlighting}[]
\FunctionTok{session\_start}\NormalTok{([}
    \StringTok{\textquotesingle{}cookie\_httponly\textquotesingle{}}\NormalTok{ =\textgreater{} }\KeywordTok{true}\OtherTok{,}
    \StringTok{\textquotesingle{}cookie\_samesite\textquotesingle{}}\NormalTok{ =\textgreater{} }\StringTok{\textquotesingle{}Strict\textquotesingle{}}\OtherTok{,}
\NormalTok{])}\OtherTok{;}
\end{Highlighting}
\end{Shaded}

\begin{itemize}
\tightlist
\item
  \texttt{cookie\_httponly} --- JavaScript cannot read the
  session cookie. This blocks XSS attacks that try to steal
  session IDs.
\item
  \texttt{cookie\_samesite} --- the browser will not send the
  cookie with requests from other sites. This is a second layer
  of CSRF defense.
\end{itemize}

\hypertarget{security-headers}{%
\section{Security headers}\label{security-headers}}

HTTP headers tell the browser how to behave. Without them, the
browser uses permissive defaults that attackers can exploit.

Add these headers in your front controller or in Apache
configuration:

\begin{Shaded}
\begin{Highlighting}[]
\CommentTok{// Prevent MIME{-}type sniffing — browser trusts the Content{-}Type header}
\FunctionTok{header}\NormalTok{(}\StringTok{\textquotesingle{}X{-}Content{-}Type{-}Options: nosniff\textquotesingle{}}\NormalTok{)}\OtherTok{;}

\CommentTok{// Deny embedding in iframes — blocks clickjacking attacks}
\FunctionTok{header}\NormalTok{(}\StringTok{\textquotesingle{}X{-}Frame{-}Options: DENY\textquotesingle{}}\NormalTok{)}\OtherTok{;}

\CommentTok{// Only send the origin as referrer, not the full URL with parameters}
\FunctionTok{header}\NormalTok{(}\StringTok{\textquotesingle{}Referrer{-}Policy: strict{-}origin{-}when{-}cross{-}origin\textquotesingle{}}\NormalTok{)}\OtherTok{;}

\CommentTok{// Control what browser features the page can use}
\FunctionTok{header}\NormalTok{(}\StringTok{\textquotesingle{}Permissions{-}Policy: camera=(), microphone=(), geolocation=()\textquotesingle{}}\NormalTok{)}\OtherTok{;}
\end{Highlighting}
\end{Shaded}

\hypertarget{content-security-policy}{%
\subsection{Content-Security-Policy}\label{content-security-policy}}

CSP is the most powerful security header. It tells the browser
exactly which resources are allowed to load:

\begin{Shaded}
\begin{Highlighting}[]
\FunctionTok{header}\NormalTok{(}\StringTok{"Content{-}Security{-}Policy: default{-}src \textquotesingle{}self\textquotesingle{}; "}
     \OperatorTok{.} \StringTok{"style{-}src \textquotesingle{}self\textquotesingle{} https://cdnjs.cloudflare.com; "}
     \OperatorTok{.} \StringTok{"font{-}src https://cdnjs.cloudflare.com; "}
     \OperatorTok{.} \StringTok{"script{-}src \textquotesingle{}self\textquotesingle{}"}\NormalTok{)}\OtherTok{;}
\end{Highlighting}
\end{Shaded}

What this means:

\begin{itemize}
\tightlist
\item
  \texttt{default-src\ \textquotesingle{}self\textquotesingle{}}
  --- by default, only load resources from our own domain.
\item
  \texttt{style-src\ \textquotesingle{}self\textquotesingle{}\ https://cdnjs.cloudflare.com}
  --- CSS from us and from the Font Awesome CDN.
\item
  \texttt{font-src\ https://cdnjs.cloudflare.com} --- fonts only
  from Font Awesome.
\item
  \texttt{script-src\ \textquotesingle{}self\textquotesingle{}}
  --- JavaScript only from our own domain. No inline scripts, no
  external scripts. If an attacker injects a
  \texttt{\textless{}script\textgreater{}} tag, the browser
  refuses to execute it.
\end{itemize}

CSP does not replace \texttt{htmlspecialchars()}. It is a safety
net --- if your escaping has a gap, CSP catches the fall.

\hypertarget{rate-limiting}{%
\section{Rate limiting}\label{rate-limiting}}

Without rate limiting, anyone can hammer your API with thousands
of requests per second. This can exhaust your server or
brute-force data through your endpoints.

A simple approach for a small application: count requests per IP
in the session.

\begin{Shaded}
\begin{Highlighting}[]
\KeywordTok{function}\NormalTok{ checkRateLimit(}\DataTypeTok{int} \VariableTok{$maxRequests} \OperatorTok{=} \DecValTok{60}\OtherTok{,} \DataTypeTok{int} \VariableTok{$windowSeconds} \OperatorTok{=} \DecValTok{60}\NormalTok{)}\OtherTok{:} \DataTypeTok{void}
\NormalTok{\{}
    \VariableTok{$now} \OperatorTok{=} \FunctionTok{time}\NormalTok{()}\OtherTok{;}
    \VariableTok{$key} \OperatorTok{=} \StringTok{\textquotesingle{}rate\_limit\textquotesingle{}}\OtherTok{;}

    \ControlFlowTok{if}\NormalTok{ (}\OperatorTok{!}\KeywordTok{isset}\NormalTok{(}\VariableTok{$\_SESSION}\NormalTok{[}\VariableTok{$key}\NormalTok{])) \{}
        \VariableTok{$\_SESSION}\NormalTok{[}\VariableTok{$key}\NormalTok{] }\OperatorTok{=}\NormalTok{ [}\StringTok{\textquotesingle{}count\textquotesingle{}}\NormalTok{ =\textgreater{} }\DecValTok{0}\OtherTok{,} \StringTok{\textquotesingle{}start\textquotesingle{}}\NormalTok{ =\textgreater{} }\VariableTok{$now}\NormalTok{]}\OtherTok{;}
\NormalTok{    \}}

    \ControlFlowTok{if}\NormalTok{ (}\VariableTok{$now} \OperatorTok{{-}} \VariableTok{$\_SESSION}\NormalTok{[}\VariableTok{$key}\NormalTok{][}\StringTok{\textquotesingle{}start\textquotesingle{}}\NormalTok{] }\OperatorTok{\textgreater{}} \VariableTok{$windowSeconds}\NormalTok{) \{}
        \VariableTok{$\_SESSION}\NormalTok{[}\VariableTok{$key}\NormalTok{] }\OperatorTok{=}\NormalTok{ [}\StringTok{\textquotesingle{}count\textquotesingle{}}\NormalTok{ =\textgreater{} }\DecValTok{0}\OtherTok{,} \StringTok{\textquotesingle{}start\textquotesingle{}}\NormalTok{ =\textgreater{} }\VariableTok{$now}\NormalTok{]}\OtherTok{;}
\NormalTok{    \}}

    \VariableTok{$\_SESSION}\NormalTok{[}\VariableTok{$key}\NormalTok{][}\StringTok{\textquotesingle{}count\textquotesingle{}}\NormalTok{]}\OperatorTok{++}\OtherTok{;}

    \ControlFlowTok{if}\NormalTok{ (}\VariableTok{$\_SESSION}\NormalTok{[}\VariableTok{$key}\NormalTok{][}\StringTok{\textquotesingle{}count\textquotesingle{}}\NormalTok{] }\OperatorTok{\textgreater{}} \VariableTok{$maxRequests}\NormalTok{) \{}
        \FunctionTok{http\_response\_code}\NormalTok{(}\DecValTok{429}\NormalTok{)}\OtherTok{;}
        \FunctionTok{header}\NormalTok{(}\StringTok{\textquotesingle{}Retry{-}After: \textquotesingle{}} \OperatorTok{.} \VariableTok{$windowSeconds}\NormalTok{)}\OtherTok{;}
        \KeywordTok{echo} \FunctionTok{json\_encode}\NormalTok{([}\StringTok{\textquotesingle{}error\textquotesingle{}}\NormalTok{ =\textgreater{} }\StringTok{\textquotesingle{}Too many requests\textquotesingle{}}\NormalTok{])}\OtherTok{;}
        \ControlFlowTok{exit}\OtherTok{;}
\NormalTok{    \}}
\NormalTok{\}}
\end{Highlighting}
\end{Shaded}

HTTP 429 (Too Many Requests) is the correct status code. The
\texttt{Retry-After} header tells the client how long to wait.

For production systems, rate limiting belongs in Apache
(\texttt{mod\_ratelimit}), Nginx, or a reverse proxy --- not in
PHP. But understanding the mechanism at the application level is
what matters for learning.

\hypertarget{https}{%
\section{HTTPS}\label{https}}

Everything above is pointless if the connection is not encrypted.
Without HTTPS:

\begin{itemize}
\tightlist
\item
  Anyone on the network can read your session cookies.
\item
  Anyone can modify the HTML in transit and inject scripts.
\item
  Security headers can be stripped by a man-in-the-middle.
\end{itemize}

In development on localhost, HTTP is fine. In production, HTTPS
is mandatory. There are no exceptions.

\hypertarget{how-it-works}{%
\subsection{How it works}\label{how-it-works}}

HTTPS uses TLS (Transport Layer Security) to encrypt the
connection. The server holds a \textbf{certificate} --- a file
that proves its identity. When the browser connects, the server
presents the certificate, they negotiate encryption keys, and
from that point every byte is encrypted.

The certificate contains the domain name, the public key, the
issuer, and an expiry date. It is signed by a \textbf{Certificate
Authority} (CA) that the browser trusts. If the signature does
not match or the certificate is expired, the browser shows a
warning.

\hypertarget{where-it-lives}{%
\subsection{Where it lives}\label{where-it-lives}}

On a Debian/Apache server, certificates typically live in
\texttt{/etc/letsencrypt/live/yourdomain/}. Two files matter:

\begin{itemize}
\tightlist
\item
  \texttt{fullchain.pem} --- the certificate (public)
\item
  \texttt{privkey.pem} --- the private key (never share this)
\end{itemize}

Apache references them in the virtual host configuration:

\begin{Shaded}
\begin{Highlighting}[]
\NormalTok{SSLCertificateFile}\StringTok{ /etc/letsencrypt/live/yourdomain/fullchain.pem}
\NormalTok{SSLCertificateKeyFile}\StringTok{ /etc/letsencrypt/live/yourdomain/privkey.pem}
\end{Highlighting}
\end{Shaded}

\hypertarget{getting-a-certificate}{%
\subsection{Getting a certificate}\label{getting-a-certificate}}

With Let's Encrypt and \texttt{certbot}, it takes two commands:

\begin{Shaded}
\begin{Highlighting}[]
\FunctionTok{sudo}\NormalTok{ apt install certbot python3{-}certbot{-}apache}
\FunctionTok{sudo}\NormalTok{ certbot }\AttributeTok{{-}{-}apache}
\end{Highlighting}
\end{Shaded}

\texttt{certbot} obtains the certificate, configures Apache, and
sets up automatic renewal. The certificate expires every 90 days,
but \texttt{certbot} renews it automatically via a systemd timer.

\hypertarget{the-layers}{%
\section{The layers}\label{the-layers}}

Security is not one thing. It is layers, each catching what the
previous one missed:

\begin{itemize}
\tightlist
\item
  \textbf{Prepared statements} → SQL injection
\item
  \textbf{Output escaping} → XSS
\item
  \textbf{CSRF tokens} → cross-site request forgery
\item
  \textbf{Input validation} → bad data, out-of-range values
\item
  \textbf{Session management} → session fixation, session theft
\item
  \textbf{Security headers} → clickjacking, MIME sniffing, CSP
  bypass
\item
  \textbf{Rate limiting} → brute force, denial of service
\item
  \textbf{HTTPS} → eavesdropping, tampering
\end{itemize}

No single layer is sufficient. Together, they make an application
that is genuinely hard to attack.

\hypertarget{the-rule}{%
\section{The rule}\label{the-rule}}

Security is not a feature you add at the end. It is a property of
every line of code you write, every header you set, every default
you accept or override.

Remember the principle from Chapter 07: each layer validates as
if it were the only one. The same applies here. Prepared
statements do not assume CSP will block the injection. CSRF
tokens do not assume HTTPS will stop the attacker. Each defense
stands on its own.

The attacker only needs one gap. You need them all closed.

\newpage

\appendix

\hypertarget{how-it-works-1}{%
\chapter{How It Works}\label{how-it-works-1}}

\hypertarget{single-point-of-control}{%
\section{Single point of control}\label{single-point-of-control}}

The application follows one rule: \textbf{the list reads, the
form writes}.

\begin{itemize}
\tightlist
\item
  \textbf{List page} --- displays all records. The only action is
  Edit.
\item
  \textbf{Form page} --- all modifications happen here: Create,
  Update, Disable, Delete.
\end{itemize}

One entity, one form, every operation. The user never has to
wonder \emph{where} to do something --- the answer is always:
open the record.

\hypertarget{the-life-of-a-request}{%
\section{The life of a request}\label{the-life-of-a-request}}

What happens when you click \textbf{Publishers} in the navigation
bar? Follow the numbers.

\begin{verbatim}
    Browser                                          Server
    ───────                                          ──────

 1  Click "Publishers"
    GET /bibliotheca/public/publishers
                │
                ▼
 2  .htaccess ─── RewriteRule ──► index.php?route=publishers
                                        │
                                        │  route in whitelist?
                                        │  yes ──► require pages/publishers.php
                                        │  no  ──► require pages/404.php
                                        │
                                        ▼
 3  Browser receives HTML ◄──────── index.php renders the full page:
    ┌─────────────────────┐         <header>, <nav>, <main>,
    │ header / nav        │         and inside <main>:
    │ ┌─────────────────┐ │         pages/publishers.php
    │ │ <table>         │ │         (empty <tbody>, plus a <script> tag)
    │ │   <thead>       │ │
    │ │   <tbody/>      │ │
    │ └─────────────────┘ │
    │ footer              │
    └─────────────────────┘
                │
                │  <script src="js/publishers.js">
                ▼
 4  PublishersView fires
    │
    │  constructor() → this.load()
    │
    │  async load() {
    │      fetch('/api/publishers.php')  ──────────────────────┐
    │  }                                                       │
    │                                                          ▼
    │                                               5  api/publishers.php
    │                                                  │
    │                                                  │  $db = new DBMS(...)
    │                                                  │  $publisher = new Publisher($db)
    │                                                  │
    │                                                  │  GET → $publisher->getAll()
    │                                                  │
    │                                                  ▼
    │                                               6  src/Publisher.php
    │                                                  │
    │                                                  │  SELECT publisher_id, name, status
    │                                                  │  FROM publisher
    │                                                  │  ORDER BY name
    │                                                  │
    │                                                  ▼
    │                                               7  src/DBMS.php
    │                                                  │
    │                                                  │  prepare → execute → fetchAll
    │                                                  │
    │                                                  ▼
    │                                               8  sql/bibliotheca.db
    │                                                  │
    │                                                  │  SQLite returns rows
    │                                                  │
    │                              JSON ◄──────────────┘
    │                              [{"publisher_id":1,"name":"Adelphi","status":1},
    │                               {"publisher_id":2,"name":"Einaudi","status":0},
    │                               ...]
    │
    │  render(publishers)
    │  │
    │  │  for each publisher:
    │  │      createElement('tr')
    │  │      if status === 0 → row.className = 'row-disabled'
    │  │      nameCell.textContent = publisher.name
    │  │      append [Edit] button
    │  │      append to <tbody>
    │  │
    ▼
 9  Browser shows the table
    ┌──────────────────────────┐
    │ header / nav             │
    │ ┌──────────────────────┐ │
    │ │ Adelphi        [Edit]│ │   ← active (normal)
    │ │ Einaudi        [Edit]│ │   ← disabled (gray)
    │ │ ...                  │ │
    │ └──────────────────────┘ │
    │ footer                   │
    └──────────────────────────┘
\end{verbatim}

\hypertarget{step-by-step}{%
\section{Step by step}\label{step-by-step}}

\hypertarget{the-click}{%
\subsection{1 --- The click}\label{the-click}}

The user clicks a link. The browser sends a GET request to
\texttt{/bibliotheca/public/publishers}.

\hypertarget{url-rewriting}{%
\subsection{2 --- URL rewriting}\label{url-rewriting}}

Apache's \texttt{.htaccess} intercepts the request. Since
\texttt{/publishers} is not a real file or directory, the
\texttt{RewriteRule} rewrites it to:

\begin{verbatim}
index.php?route=publishers
\end{verbatim}

\hypertarget{the-router}{%
\subsection{3 --- The router}\label{the-router}}

\texttt{index.php} reads the \texttt{route} parameter, checks it
against a whitelist of allowed routes, and includes the matching
page inside a common HTML shell (header, nav, main, footer).

The page \texttt{pages/publishers.php} is minimal: an empty
\texttt{\textless{}table\textgreater{}} and a
\texttt{\textless{}script\textgreater{}} tag. No data yet.

\hypertarget{javascript-takes-over}{%
\subsection{4 --- JavaScript takes
over}\label{javascript-takes-over}}

The browser loads \texttt{js/publishers.js}. The script creates a
\texttt{PublishersView} object. The constructor immediately calls
\texttt{load()}, which uses \texttt{fetch} to request data from
the API.

\hypertarget{the-controller-1}{%
\subsection{5 --- The controller}\label{the-controller-1}}

\texttt{api/publishers.php} receives the GET request. It creates
a \texttt{DBMS} instance (the database connection) and a
\texttt{Publisher} instance (the model). Based on the HTTP
method, it calls the right model method.

For a plain GET with no \texttt{id} parameter, it calls
\texttt{getAll()}.

\hypertarget{the-model-1}{%
\subsection{6 --- The model}\label{the-model-1}}

\texttt{Publisher::getAll()} builds a SQL query that selects all
records --- both active and disabled --- and delegates execution
to DBMS.

The model knows \emph{what} to ask, but not \emph{how} the
database works. That is the DBMS wrapper's job.

\hypertarget{the-database-wrapper}{%
\subsection{7 --- The database
wrapper}\label{the-database-wrapper}}

\texttt{DBMS::fetchAll()} prepares the statement, executes it,
and returns an array of associative arrays. Every query goes
through prepared statements --- never string concatenation.

\hypertarget{sqlite-1}{%
\subsection{8 --- SQLite}\label{sqlite-1}}

The database engine runs the query against
\texttt{sql/bibliotheca.db} and returns the rows to PHP.

\hypertarget{rendering}{%
\subsection{9 --- Rendering}\label{rendering}}

The JSON response travels back to the browser.
\texttt{PublishersView.render()} loops through the array, creates
a \texttt{\textless{}tr\textgreater{}} for each publisher using
\texttt{createElement} and \texttt{textContent} (never
\texttt{innerHTML}), and appends everything to the
\texttt{\textless{}tbody\textgreater{}}.

Rows with \texttt{status\ ===\ 0} get the class
\texttt{row-disabled} and appear grayed out. The list is
read-only --- the only action is the Edit button, which links to
the form page.

\hypertarget{the-round-trip}{%
\section{The round trip}\label{the-round-trip}}

The key insight: \textbf{the page loads twice}.

\begin{enumerate}
\def\labelenumi{\arabic{enumi}.}
\tightlist
\item
  \textbf{First trip} --- the browser gets the HTML shell (step
  1--3). The table is empty.
\item
  \textbf{Second trip} --- JavaScript fetches the data as JSON
  (step 4--9). The table fills up.
\end{enumerate}

This is the foundation of every modern web application. The
structure (HTML) and the data (JSON) travel separately. The
browser assembles them.

\hypertarget{what-travels-on-the-wire}{%
\section{What travels on the
wire}\label{what-travels-on-the-wire}}

When the browser sends a request, it is not a file. It is a block
of text transmitted over a TCP connection to port 80 (or 443 for
HTTPS). You can see it with \texttt{curl\ -v}:

\begin{verbatim}
> GET /bibliotheca/public/categories HTTP/1.1
> Host: localhost
> User-Agent: curl/7.88.1
> Accept: */*
>
\end{verbatim}

The first line is the \textbf{request line}: method, path,
protocol. The lines that follow are \textbf{headers}: key-value
pairs that describe who is calling, what formats are accepted,
the session cookie, and so on. The empty line marks the end. A
GET request has no body; a POST would carry one after the blank
line.

The server responds in the same format: a status line, headers, a
blank line, then the body (HTML, JSON, or whatever was asked).

\begin{verbatim}
< HTTP/1.1 200 OK
< Content-Type: text/html; charset=UTF-8
<
< <!DOCTYPE html>...
\end{verbatim}

No files are written to disk. Apache has a process (managed by
its \textbf{MPM}, Multi-Processing Module) listening on a TCP
socket. The bytes arrive in memory, Apache interprets them, calls
PHP, and sends the response back on the same connection.

With HTTP/1.1, the connection stays open briefly
(\textbf{keep-alive}) so the browser can reuse it for subsequent
requests (the CSS file, the JavaScript, the favicon) without
opening a new one.

\hypertarget{the-dom}{%
\section{The DOM}\label{the-dom}}

When the browser receives HTML, it does not keep it as text. It
parses it and builds an in-memory tree of objects called the
\textbf{DOM} (Document Object Model). For example:

\begin{Shaded}
\begin{Highlighting}[]
\KeywordTok{\textless{}table} \ErrorTok{id}\OtherTok{=}\StringTok{"category{-}table"}\KeywordTok{\textgreater{}}
    \KeywordTok{\textless{}thead\textgreater{}}
        \KeywordTok{\textless{}tr\textgreater{}\textless{}th\textgreater{}}\NormalTok{Name}\KeywordTok{\textless{}/th\textgreater{}\textless{}/tr\textgreater{}}
    \KeywordTok{\textless{}/thead\textgreater{}}
    \KeywordTok{\textless{}tbody\textgreater{}\textless{}/tbody\textgreater{}}
\KeywordTok{\textless{}/table\textgreater{}}
\end{Highlighting}
\end{Shaded}

becomes:

\begin{verbatim}
table#category-table
├── thead
│   └── tr
│       └── th → "Name"
└── tbody (empty)
\end{verbatim}

Every tag is a \textbf{node} with properties and methods.
JavaScript works on this tree, never on the HTML source. When it
calls
\texttt{document.createElement(\textquotesingle{}tr\textquotesingle{})}
and \texttt{tbody.appendChild(row)}, it is adding nodes to the
tree. The browser detects the change and updates the screen.

This is why the \texttt{\textless{}script\textgreater{}} tag sits
at the bottom of the page template: by that point the browser has
already built the DOM, and JavaScript can find the elements it
needs with
\texttt{document.querySelector(\textquotesingle{}\#category-table\ tbody\textquotesingle{})}.

The selector syntax (\texttt{\#id}, \texttt{.class},
\texttt{tag}) is the same one used in CSS. JavaScript borrows it
to navigate the tree.

\hypertarget{how-data-changes-shape}{%
\section{How data changes shape}\label{how-data-changes-shape}}

A single piece of data (say, a category name) passes through four
representations on its way from disk to screen:

\begin{longtable}[]{@{}ll@{}}
\toprule
Layer & Form \\
\midrule
\endhead
SQLite & Row in a table \\
PHP & Associative array \\
Network & JSON text \\
JavaScript & Object with properties \\
\bottomrule
\end{longtable}

\begin{verbatim}
SQLite:  | 3 | Computer Science | 1 |
              ↓
PHP:     ['category_id'=>3,
          'name'=>'Computer Science',
          'status'=>1]
              ↓
JSON:    {"category_id":3,
          "name":"Computer Science",
          "status":1}
              ↓
JS:      {category_id: 3,
          name: "Computer Science",
          status: 1}
\end{verbatim}

Each layer only knows about the one directly below it. JavaScript
has no idea SQLite exists; it just sees a URL that returns JSON.
The model has no idea JavaScript exists; it just returns an array
to the controller.

\hypertarget{where-each-file-lives}{%
\section{Where each file lives}\label{where-each-file-lives}}

\begin{verbatim}
bibliotheca/
    src/                    ← private (not web-accessible)
        DBMS.php            7  database wrapper
        Publisher.php       6  model
    sql/
        bibliotheca.db      8  database
    public/                 ← document root (web-accessible)
        .htaccess           2  URL rewriting
        index.php           3  router
        pages/
            publishers.php  3  list template
            publisher.php      form template
        js/
            publishers.js   4  list logic
            publisher.js       form logic
        api/
            publishers.php  5  controller
\end{verbatim}

The numbers match the steps in the diagram. Notice how
\texttt{src/} and \texttt{sql/} sit outside \texttt{public/}. The
browser can never reach them directly --- only PHP can.

\hypertarget{the-life-of-a-form}{%
\section{The life of a form}\label{the-life-of-a-form}}

The list page shows data. The form page changes it. What happens
when you click \textbf{Add new publisher}, or \textbf{Edit} on an
existing one?

\hypertarget{opening-the-form}{%
\subsection{Opening the form}\label{opening-the-form}}

\begin{verbatim}
    Browser                                          Server
    ───────                                          ──────

    ┌─────────────────────────────────────┐
    │  Two ways to reach the form:        │
    │                                     │
    │  A) Click "Add new publisher"       │
    │     GET /publisher                  │
    │     (no ?id parameter)              │
    │                                     │
    │  B) Click "Edit" on Einaudi         │
    │     GET /publisher?id=2             │
    │                                     │
    └─────────────────────────────────────┘
                │
                ▼
 1  .htaccess ─── RewriteRule ──► index.php?route=publisher
                                        │
                                        ▼
 2  Browser receives HTML ◄──────── index.php renders the page:
    ┌─────────────────────┐         pages/publisher.php
    │ header / nav        │
    │ ┌─────────────────┐ │         <form> with:
    │ │ Add Publisher    │ │         - hidden <input> for ID (empty)
    │ │                  │ │         - text <input> for name (empty)
    │ │ Name: [        ] │ │         - checkbox Active (hidden until Edit)
    │ │                  │ │         - Save button
    │ │ [Save] [Cancel]  │ │         - Delete button (hidden until Edit)
    │ └─────────────────┘ │
    │ footer              │
    └─────────────────────┘
                │
                │  <script src="js/publisher.js">
                ▼
 3  PublisherForm fires
    │
    │  constructor()
    │  │  binds form submit → this.save()
    │  │  binds delete button → this.remove()
    │  │  calls this.checkEdit()
    │  │
    │  checkEdit()
    │  │  reads ?id from URL
    │  │
    │  ├─── no ?id ──► form stays empty ("Add Publisher")
    │  │                checkbox and Delete stay hidden
    │  │                done — waiting for user input
    │  │
    │  └─── ?id=2 ──► this.load(2)
    │                   │
    │                   │  fetch('/api/publishers.php?id=2')  ─────┐
    │                   │                                          │
    │                   │                                          ▼
    │                   │                               api/publishers.php
    │                   │                               │
    │                   │                               │  GET with ?id=2
    │                   │                               │  $publisher->getById(2)
    │                   │                               │
    │                   │                               │  SELECT ... WHERE publisher_id = :id
    │                   │                               │
    │                   │  JSON ◄────────────────────────┘
    │                   │  {"publisher_id":2,"name":"Einaudi","status":1}
    │                   │
    │                   │  fills the form:
    │                   │  - hidden ID ← 2
    │                   │  - name input ← "Einaudi"
    │                   │  - checkbox Active ← checked (status 1)
    │                   │  - shows checkbox and Delete button
    │                   │  - title ← "Edit Publisher"
    │                   │
    │                   ▼
    │  ┌───────────────────────┐
    │  │ Edit Publisher         │
    │  │                        │
    │  │ Name: [Einaudi      ]  │
    │  │ ☑ Active               │
    │  │                        │
    │  │ [Save] [Delete] Cancel │
    │  └────────────────────────┘
    │
    ▼  Waiting for user input.
\end{verbatim}

\hypertarget{saving-the-form}{%
\subsection{Saving the form}\label{saving-the-form}}

\begin{verbatim}
    Browser                                          Server
    ───────                                          ──────

 4  User modifies the record, clicks Save
    │
    │  form submit event
    │  e.preventDefault() — no page reload
    │
    │  this.save()
    │  │
    │  │  this.validate()
    │  │  ├─── name empty? ──► show error, stop
    │  │  └─── name ok?    ──► continue
    │  │
    │  │  check hidden ID field:
    │  │
    │  ├─── ID is empty (Add mode)
    │  │    │
    │  │    │  fetch(API, {                    ────────────────────┐
    │  │    │      method: 'POST',                                │
    │  │    │      body: {"name":"Einaudi Editore"}               │
    │  │    │  })                                                  │
    │  │    │                                                      ▼
    │  │    │                                           api/publishers.php
    │  │    │                                           │
    │  │    │                                           │  POST
    │  │    │                                           │  validate: name required? ✓
    │  │    │                                           │  check: exists()? → no
    │  │    │                                           │  $publisher->insert("Einaudi Editore")
    │  │    │                                           │
    │  │    │                                           │  INSERT INTO publisher (name)
    │  │    │                                           │  VALUES (:name)
    │  │    │                                           │
    │  │    │  JSON ◄───────────────────────────────────┘
    │  │    │  {"publisher_id":13,"name":"Einaudi Editore"}
    │  │
    │  └─── ID is 2 (Edit mode)
    │       │
    │       │  fetch(API, {                    ────────────────────┐
    │       │      method: 'PUT',                                  │
    │       │      body: {"publisher_id":2,                        │
    │       │             "name":"Einaudi Editore",                │
    │       │             "status":0}                              │
    │       │  })                                                   │
    │       │                                                       ▼
    │       │                                            api/publishers.php
    │       │                                            │
    │       │                                            │  PUT
    │       │                                            │  validate: ID and name? ✓
    │       │                                            │  check: exists()? → no
    │       │                                            │  status=0 and hasBooks()? → block or allow
    │       │                                            │  $publisher->update(2, "Einaudi Editore", 0)
    │       │                                            │
    │       │                                            │  UPDATE publisher
    │       │                                            │  SET name = :name, status = :status
    │       │                                            │  WHERE publisher_id = :id
    │       │                                            │
    │       │  JSON ◄────────────────────────────────────┘
    │       │  {"publisher_id":2,"name":"Einaudi Editore","status":0}
    │
    │
 5  ├─── response ok? ──► redirect to /publishers (the list reloads)
    │
    └─── response error (409/422)?
         │
         │  {"error":"Publisher already exists"}
         │  — or —
         │  {"error":"Cannot disable: publisher has associated books"}
         │
         │  this.showError('publisher-name', result.error)
         │
         ▼
         ┌──────────────────────────────┐
         │ Edit Publisher                │
         │                               │
         │ Name: [Einaudi Editore     ]  │
         │ Publisher already exists       │
         │ ☐ Active                      │
         │                               │
         │ [Save] [Delete] Cancel        │
         └───────────────────────────────┘
\end{verbatim}

\hypertarget{deleting-from-the-form}{%
\subsection{Deleting from the
form}\label{deleting-from-the-form}}

\begin{verbatim}
    Browser                                          Server
    ───────                                          ──────

    ┌───────────────────────┐
    │ Edit Publisher         │
    │                        │
    │ Name: [Einaudi      ]  │
    │ ☑ Active               │
    │                        │
    │ [Save] [Delete] Cancel │  ◄── user clicks [Delete]
    └────────────────────────┘
                │
                ▼
 1  confirm('Permanently delete this publisher?
             This cannot be undone.')
    │
    ├─── Cancel ──► nothing happens, form stays open
    │
    └─── OK
         │
         │  fetch(API, {                         ────────────────────┐
         │      method: 'DELETE',                                    │
         │      body: {"publisher_id": 2}                            │
         │  })                                                       │
         │                                                           ▼
         │                                                api/publishers.php
         │                                                │
         │                                                │  DELETE
         │                                                │  validate: ID required? ✓
         │                                                │
         │                                                │  $publisher->hasBooks(2, false)?
         │                                                │  (checks ALL books, active or not)
         │                                                │
         │                                                ├─── yes (has books)
         │                                                │    │
         │                                                │    │  HTTP 422
         │                                                │    │  {"error":"Cannot delete:
         │                                                │    │   publisher has associated books"}
         │                                                │    │
         │                                                └─── no (safe to delete)
         │                                                     │
         │                                                     │  $publisher->delete(2)
         │                                                     │
         │                                                     │  DELETE FROM publisher
         │                                                     │  WHERE publisher_id = :id
         │                                                     │
         │                                                     │  {"deleted": true}
         │                                                     │
 2  JSON ◄─────────────────────────────────────────────────────┘
    │
    ├─── response ok? ──► redirect to /publishers
    │
    └─── response error (409/422)?
         │
         │  alert("Cannot delete: publisher has associated books")
         │  form stays open
         ▼
\end{verbatim}

\hypertarget{step-by-step-1}{%
\subsection{Step by step}\label{step-by-step-1}}

\textbf{Step 1--2} --- identical to the list page:
\texttt{.htaccess} rewrites, the router loads
\texttt{pages/publisher.php} inside the HTML shell. The form
arrives empty.

\textbf{Step 3 --- checkEdit} --- \texttt{PublisherForm} reads
the URL. If there is no \texttt{?id}, the form stays empty --- it
is an Add. The Active checkbox and Delete button stay hidden
because they make no sense on a record that does not exist yet.

If \texttt{?id=2} is present, JavaScript fetches that single
record from the API and fills the form fields, including the
checkbox state. The title changes from ``Add'' to ``Edit'', and
the checkbox and Delete button become visible.

One form, two modes. The hidden
\texttt{\textless{}input\textgreater{}} for the ID decides which.

\textbf{Step 4 --- save} --- the user clicks Save. JavaScript
prevents the default form submission (no page reload) and runs
validation. If the name is empty, it shows an inline error and
stops.

If validation passes, JavaScript checks the hidden ID: -
\textbf{Empty} → POST (create a new record) - \textbf{Has a
value} → PUT (update the existing record, including status)

The server validates again (never trust the client), normalizes
the input (\texttt{ucwords}, \texttt{strtolower}, \texttt{trim}),
checks for duplicates, and checks whether a disable is allowed (a
publisher with active books cannot be disabled). Then it executes
the query.

\textbf{Step 5 --- response} --- if the server returns OK,
JavaScript redirects to the list page. If the server returns an
error (e.g., 409 Conflict for a duplicate name, or 422 for
disabling with active books), JavaScript shows the server's error
message inline under the field.

\textbf{Delete} --- happens from the same form page. The user
clicks Delete, confirms the dialog, and JavaScript sends a DELETE
request. The server checks referential integrity --- a publisher
with any books (active or disabled) cannot be deleted. On
success, the record is permanently removed
(\texttt{DELETE\ FROM}) and the user is redirected to the list.

\hypertarget{disable-vs-delete}{%
\subsection{Disable vs Delete}\label{disable-vs-delete}}

The form offers two ways to remove a record:

\begin{itemize}
\tightlist
\item
  \textbf{Disable} --- uncheck Active, click Save. SQL:
  \texttt{UPDATE\ status=0}. Reversible. Checks for active books.
\item
  \textbf{Delete} --- click Delete. SQL: \texttt{DELETE\ FROM}.
  Permanent. Checks for any books.
\end{itemize}

Disabling keeps the record in the database but grayed out in the
list. It can be reversed by checking Active again.

Deleting removes the record forever. The safety check is
stricter: it looks for \emph{any} associated books, not just
active ones, because even a disabled book still holds a foreign
key reference.

\hypertarget{the-key-insight}{%
\subsection{The key insight}\label{the-key-insight}}

The form page loads \textbf{once or twice}, depending on the
mode:

\begin{longtable}[]{@{}lll@{}}
\toprule
Mode & Trips & What happens \\
\midrule
\endhead
Add & 1 & HTML arrives, form is empty, done \\
Edit & 2 & HTML arrives, then fetch loads the record \\
\bottomrule
\end{longtable}

Saving always adds \textbf{one more trip}: the POST or PUT to the
API. On success, the redirect to the list triggers the full list
flow again (steps 1--9 from the first diagram).

\hypertarget{the-complete-picture}{%
\section{The complete picture}\label{the-complete-picture}}

All four CRUD operations, one entity, two pages:

\begin{itemize}
\tightlist
\item
  \textbf{Read} --- list page, \texttt{GET}, 2 trips. Table fills
  with data.
\item
  \textbf{Create} --- form page, \texttt{POST}, 2 trips. Empty
  form → save → list.
\item
  \textbf{Update} --- form page, \texttt{PUT}, 3 trips. Form
  loads record → save → list.
\item
  \textbf{Disable} --- form page, \texttt{PUT}, 3 trips. Uncheck
  Active → save → gray row.
\item
  \textbf{Delete} --- form page, \texttt{DELETE}, 2 trips. Click
  Delete → confirm → gone.
\end{itemize}

The list is read-only. Every write operation goes through the
form. This is the single point of control.

Every other entity (Category, Author, Book) follows the same
pattern. The files change, the flow does not.

\hypertarget{authentication}{%
\section{Authentication}\label{authentication}}

Bibliotheca has a single admin user. The login flow adds one
layer on top of everything described above.

\hypertarget{the-login}{%
\subsection{The login}\label{the-login}}

\begin{enumerate}
\def\labelenumi{\arabic{enumi}.}
\tightlist
\item
  The user visits \texttt{/login} (default credentials:
  \texttt{admin} / \texttt{bibliotheca}). The front controller
  loads \texttt{pages/login.php}, which includes
  \texttt{js/login.js}.
\item
  The user submits username and password. JavaScript sends a POST
  to \texttt{api/auth.php} with JSON body.
\item
  The API calls \texttt{Auth::login()}, which queries the
  \texttt{user} table and compares the password with
  \texttt{password\_verify()}.
\item
  If valid, \texttt{session\_regenerate\_id(true)} creates a new
  session (preventing fixation), and
  \texttt{\$\_SESSION{[}\textquotesingle{}user\_id\textquotesingle{}{]}}
  is set. The API returns \texttt{\{"authenticated":\ true\}}.
\item
  JavaScript redirects to the home page.
\end{enumerate}

\hypertarget{changing-the-password}{%
\subsection{Changing the password}\label{changing-the-password}}

The default password should be changed immediately. Generate a
new hash from the command line:

\begin{Shaded}
\begin{Highlighting}[]
\ExtensionTok{php} \AttributeTok{{-}r} \StringTok{"echo password\_hash(\textquotesingle{}yournewpassword\textquotesingle{}, PASSWORD\_DEFAULT);"}
\end{Highlighting}
\end{Shaded}

Then update the database:

\begin{Shaded}
\begin{Highlighting}[]
\ExtensionTok{sqlite3}\NormalTok{ sql/bibliotheca.db}
\ExtensionTok{UPDATE}\NormalTok{ user SET password = }\StringTok{\textquotesingle{}$2y$10$...\textquotesingle{}}\NormalTok{ WHERE username = }\StringTok{\textquotesingle{}admin\textquotesingle{}}\KeywordTok{;}
\end{Highlighting}
\end{Shaded}

Replace \texttt{\$2y\$10\$...} with the hash you generated. The
username can be changed the same way.

\hypertarget{protecting-writes}{%
\subsection{Protecting writes}\label{protecting-writes}}

Every API endpoint checks authentication before processing POST,
PUT, or DELETE requests:

\begin{Shaded}
\begin{Highlighting}[]
\ConstantTok{C}\NormalTok{srf::start()}\OtherTok{;}
\ConstantTok{C}\NormalTok{srf::verify()}\OtherTok{;}
\ConstantTok{A}\NormalTok{uth::}\KeywordTok{require}\NormalTok{()}\OtherTok{;}  \CommentTok{// 401 if not logged in}
\end{Highlighting}
\end{Shaded}

\texttt{Auth::require()} reads
\texttt{\$\_SESSION{[}\textquotesingle{}user\_id\textquotesingle{}{]}}.
If it is not set, it responds with HTTP 401 and exits.

\hypertarget{hiding-the-ui}{%
\subsection{Hiding the UI}\label{hiding-the-ui}}

The front controller sets \texttt{\$isLoggedIn\ =\ Auth::check()}
and exposes it to templates and JavaScript:

\begin{itemize}
\tightlist
\item
  \textbf{PHP templates} --- the ``Add new'' buttons are wrapped
  in
  \texttt{\textless{}?php\ if\ (\$isLoggedIn):\ ?\textgreater{}}
  blocks.
\item
  \textbf{JavaScript} --- the
  \texttt{\textless{}meta\ name="authenticated"\textgreater{}}
  tag carries the state. \texttt{auth.js} reads it into
  \texttt{AUTH.authenticated}, and list views check this before
  rendering Edit buttons.
\item
  \textbf{Form pages} --- the front controller redirects to
  \texttt{/login} if the user tries to access a form page without
  being logged in.
\end{itemize}

The API is the real guard. The UI changes are convenience ---
they prevent confusion, not attacks.

\hypertarget{see-also}{%
\section{See also}\label{see-also}}

This overview shows the full journey. The chapters break it into
layers:

\begin{itemize}
\tightlist
\item
  \href{03_structure.md}{Chapter 03 --- Project Structure} ---
  where files live and why
\item
  \href{04_backend.md}{Chapter 04 --- Backend} --- DBMS, Model,
  Controller
\item
  \href{05_frontend.md}{Chapter 05 --- Frontend} --- routing,
  fetch, DOM
\item
  \href{06_crud.md}{Chapter 06 --- CRUD} --- the four operations
  in detail
\end{itemize}

\newpage

\hypertarget{glossary}{%
\chapter{Glossary}\label{glossary}}

Terms you will encounter in this book, in alphabetical order.

\begin{center}\rule{0.5\linewidth}{0.5pt}\end{center}

\textbf{42} --- The answer to the ultimate question of life, the
universe, and everything. From Douglas Adams' \emph{The
Hitchhiker's Guide to the Galaxy} (1979). The joke is that nobody
knows the question. In programming, it is a reminder that knowing
the answer is useless if you do not understand the problem.

\textbf{\texttt{.htaccess}} --- A configuration file read by
Apache on every request to the directory where it lives. The dot
prefix makes it hidden on Unix systems. In Bibliotheca,
\texttt{.htaccess} contains a single rewrite rule: any URL that
does not match an existing file is forwarded to
\texttt{index.php?route=...}, enabling clean URLs like
\texttt{/publishers} instead of
\texttt{index.php?route=publishers}. Requires Apache's
\texttt{mod\_rewrite} module to be enabled.

\textbf{API (Application Programming Interface)} --- A set of
endpoints that a program can call to exchange data. In
Bibliotheca, the files in \texttt{public/api/} are the API:
JavaScript sends HTTP requests, PHP answers with JSON.

\textbf{async/await} --- JavaScript syntax that makes
asynchronous code read like sequential code.
\texttt{await\ fetch(url)} pauses execution until the response
arrives, without blocking the browser.

\textbf{Böhm-Jacopini theorem} --- Any computable function can be
written using only three control structures: sequence, selection
(if/else), and iteration (for/while). No goto, no magic.

\textbf{CRUD} --- Create, Read, Update, Delete --- the four basic
operations on data. Every entity in Bibliotheca implements all
four through a chain of HTTP method → Controller → Model → SQL.

\begin{longtable}[]{@{}lll@{}}
\toprule
Operation & HTTP & SQL \\
\midrule
\endhead
Create & POST & INSERT \\
Read & GET & SELECT \\
Update & PUT & UPDATE \\
Delete & DELETE & DELETE \\
\bottomrule
\end{longtable}

\textbf{Controller} --- The PHP script in \texttt{public/api/}
that receives an HTTP request, calls the Model, and returns JSON.
It decides \emph{what to do} but does not know \emph{how to
query} or \emph{how to display}.

\textbf{curl} --- A command-line tool for making HTTP requests.
Useful for testing APIs without a browser. In Bibliotheca,
\texttt{curl} lets you call the API endpoints directly from the
terminal:
\texttt{curl\ http://localhost/bibliotheca/public/api/publishers.php}.
The name stands for ``Client URL''.

\textbf{DDL / DML / DQL} --- Three categories of SQL: -
\textbf{DDL} (Data Definition Language) ---
\texttt{CREATE\ TABLE}, \texttt{ALTER\ TABLE} - \textbf{DML}
(Data Manipulation Language) --- \texttt{INSERT},
\texttt{UPDATE}, \texttt{DELETE} - \textbf{DQL} (Data Query
Language) --- \texttt{SELECT}

In Bibliotheca, each has its own directory under \texttt{sql/}.

\textbf{Dependency injection} --- Passing an object to a class
instead of letting the class create it internally. The Model
receives the DBMS in its constructor --- it does not know how the
connection was opened, only that it can use it. This makes the
code testable and flexible.

\textbf{Document root} --- The directory that Apache serves to
browsers. In Bibliotheca, this is \texttt{public/}. Files outside
it (\texttt{src/}, \texttt{sql/}) are invisible to the web ---
only PHP can reach them.

\textbf{DOM (Document Object Model)} --- The browser's in-memory
representation of the HTML page. JavaScript manipulates the DOM
to build tables, add rows, show errors --- all without reloading
the page.

\textbf{DRY (Don't Repeat Yourself)} --- Every piece of knowledge
should have a single representation. If you find yourself copying
the same code, extract it.

\textbf{ES6} --- The 2015 edition of JavaScript that introduced
classes, \texttt{const}/\texttt{let}, arrow functions, template
literals, and \texttt{async/await}. Bibliotheca uses ES6 classes
on the frontend.

\textbf{Fail Fast} --- If something goes wrong, stop immediately.
Do not drag an error through the code hoping it will resolve
itself. In DBMS.php, \texttt{PDO::ERRMODE\_EXCEPTION} enforces
this: a bad query throws an exception instead of returning false.

\textbf{Fail Safe} --- When you fail, fail securely. Never leave
the system in an inconsistent or unsafe state.

\textbf{\texttt{fetch}} --- The browser's native API for HTTP
requests. Replaces the older \texttt{XMLHttpRequest}. Combined
with \texttt{async/await}, it is how JavaScript talks to the
backend.

\textbf{Foreign key} --- A column that references the primary key
of another table. The \texttt{publisher\_id} column in
\texttt{book} points to \texttt{publisher.publisher\_id}. SQLite
ignores foreign keys by default ---
\texttt{PRAGMA\ foreign\_keys\ =\ ON} enables enforcement.

\textbf{Front controller} --- A design pattern where every
request enters the application through a single file. In
Bibliotheca, \texttt{index.php} is the front controller:
\texttt{.htaccess} routes all URLs to it, and it decides which
page to load.

\textbf{Hacker} --- Originally, someone who explores systems out
of curiosity, taking things apart to understand how they work.
The media turned the word into a synonym for criminal, but in
programming culture a hacker is simply a person driven by the
urge to learn by doing. That is the meaning we use in this book.

\textbf{HTTP methods} --- The verbs of the web: GET (read), POST
(create), PUT (update), DELETE (remove). Each API endpoint in
Bibliotheca uses the method to decide which operation to perform.

\textbf{HTTP status codes} --- Numbers that tell the client what
happened. The ones Bibliotheca uses: - \textbf{200} --- OK -
\textbf{400} --- Bad Request (missing or invalid input) -
\textbf{404} --- Not Found - \textbf{405} --- Method Not Allowed
- \textbf{401} --- Unauthorized (not authenticated) -
\textbf{409} --- Conflict (duplicate record) - \textbf{422} ---
Unprocessable Entity (dependency error, invalid data) -
\textbf{429} --- Too Many Requests (rate limiting) - \textbf{500}
--- Internal Server Error

\textbf{Index (database)} --- A data structure that speeds up
lookups on a column, like a book index speeds up finding a topic.
SQLite creates indexes automatically for primary keys, but
foreign key columns need explicit \texttt{CREATE\ INDEX}
statements. Without an index, the database must scan every row
(full table scan) to find a match.

\textbf{\texttt{innerHTML} (avoid)} --- A DOM property that
parses a string as HTML. If the string contains user input, this
is an XSS vulnerability. Bibliotheca uses \texttt{textContent}
instead, which treats everything as plain text.

\textbf{JSON (JavaScript Object Notation)} --- A lightweight data
format. The API returns JSON, JavaScript parses it with
\texttt{response.json()}. It looks like this:
\texttt{\{"name":\ "Einaudi",\ "publisher\_id":\ 2\}}.

\textbf{Junction table} --- A table that implements a
many-to-many relationship. \texttt{book\_author} connects books
to authors: one book can have many authors, one author can have
many books.

\textbf{kebab-case} --- Words separated by hyphens:
\texttt{book-list}, \texttt{main-header}. Used for HTML IDs, CSS
classes, and git branch names.

\textbf{KISS (Keep It Simple, Stupid)} --- Always choose the
simplest solution that works. Complexity is the enemy of
understanding.

\textbf{\texttt{mod\_rewrite}} --- An Apache module that rewrites
URLs before they reach your code. Without it, \texttt{.htaccess}
rewrite rules are silently ignored. Enable it with
\texttt{sudo\ a2enmod\ rewrite} and restart Apache. It is what
makes clean URLs possible --- turning \texttt{/publishers} into
\texttt{index.php?route=publishers} behind the scenes.

\textbf{Model} --- A PHP class in \texttt{src/} that knows how to
query one entity. It receives the DBMS via dependency injection
and provides methods like \texttt{getAll}, \texttt{getById},
\texttt{insert}, \texttt{update}, \texttt{delete}.

\textbf{MVC (Model-View-Controller)} --- An architecture that
separates data (Model), presentation (View), and logic
(Controller). In Bibliotheca: - \textbf{Model} ---
\texttt{src/*.php} - \textbf{View} ---
\texttt{public/pages/*.php} + \texttt{public/js/*.js} -
\textbf{Controller} --- \texttt{public/api/*.php}

\textbf{PDO (PHP Data Objects)} --- PHP's database abstraction
layer. It works the same way regardless of the database engine
(SQLite, MySQL, PostgreSQL). Switch database by changing the
connection string, nothing else.

\textbf{Prepared statement} --- A SQL query with placeholders
instead of values. The database engine separates the query from
the data, making SQL injection impossible. Every query in
Bibliotheca uses them.

PDO supports two styles of placeholders: - \textbf{Named} ---
\texttt{:id}, \texttt{:name}. Self-documenting,
order-independent. - \textbf{Positional} --- \texttt{?}. Shorter,
but you must match the order.

We use named parameters because they are easier to read and safer
to refactor.
\texttt{WHERE\ publisher\_id\ =\ :id\ AND\ status\ =\ :status}
tells you what each value is.
\texttt{WHERE\ publisher\_id\ =\ ?\ AND\ status\ =\ ?} requires
counting positions.

\textbf{Primary key} --- A column that uniquely identifies each
row. In Bibliotheca, it is always the first column:
\texttt{publisher\_id}, \texttt{book\_id}, etc.

\textbf{PSR-12} --- A PHP coding standard that defines formatting
rules: indentation, brace placement, naming. Bibliotheca follows
it.

\textbf{Routing} --- The technique of directing every request
through a single entry point (\texttt{index.php}) that decides
which page to load. One door, many rooms. If a route is not in
the whitelist, it gets a 404. Every modern web framework uses
this pattern.

\textbf{Rewrite rule} --- An Apache directive in
\texttt{.htaccess} that transforms URLs before they reach PHP.
The rule \texttt{\^{}(.*)\$\ index.php?route=\$1} captures the
URL path and passes it as a parameter to the router.

\textbf{REST (Representational State Transfer)} --- An
architectural style for web APIs, defined by Roy Fielding in
2000. The idea is simple: use what HTTP already gives you. -
\textbf{Resources have URLs.} Each entity lives at its own
address: \texttt{/api/books.php}, \texttt{/api/publishers.php}. -
\textbf{Verbs come from HTTP.} The method \emph{is} the action:
GET reads, POST creates, PUT updates, DELETE removes. No need for
\texttt{?action=delete} in the URL. - \textbf{Responses are
representations.} The server returns data (JSON), not pages. The
client decides how to display it. - \textbf{Each request is
self-contained.} The server does not remember previous requests.
Every call carries all the information it needs.

Before REST, web services used complex protocols like SOAP ---
XML-heavy, with action names buried inside the message body. REST
replaced all of that with what the web already had: URLs for
nouns, HTTP methods for verbs, and JSON for data. Bibliotheca
follows this style.

\textbf{SoC (Separation of Concerns)} --- Each part of the code
has one job. Backend serves data, frontend presents it. Model
queries, Controller routes, View displays.

\textbf{Soft delete vs hard delete} --- Two strategies for
removing data. Soft delete sets a flag (\texttt{status\ =\ 0})
--- the record stays in the database but appears disabled. Hard
delete runs \texttt{DELETE\ FROM} --- the record is gone.
Bibliotheca uses both: the status toggle disables a record
(greyed out in the list), while the Delete button removes it
permanently.

\textbf{SPA (Single Page Application)} --- A web application
where the browser loads one HTML page and JavaScript handles all
navigation and data updates via \texttt{fetch} calls to the API.
The page never fully reloads --- only the content changes.
Bibliotheca follows this pattern: \texttt{index.php} serves the
shell, JavaScript builds the interface.

\textbf{SQL injection} --- An attack where user input is inserted
into a SQL query, changing its meaning. Example: entering
\texttt{\textquotesingle{};\ DROP\ TABLE\ book;\ -\/-} as a name.
Prepared statements prevent this completely.

\textbf{SQLite} --- A database engine stored in a single file. No
server process, no configuration. The file
\texttt{sql/bibliotheca.db} is the entire database.

\textbf{\texttt{strict\_types}} ---
\texttt{declare(strict\_types=1);} at the top of a PHP file
enables strict type checking. Passing a string where an int is
expected throws a TypeError instead of silently converting.

\textbf{tail} --- A Unix command that shows the last lines of a
file. Essential for reading logs:
\texttt{sudo\ tail\ -20\ /var/log/apache2/error.log} shows the 20
most recent entries. Use \texttt{tail\ -f} to follow a log file
in real time as new entries appear.

\textbf{\texttt{textContent}} --- A DOM property that sets or
gets the text of an element. Unlike \texttt{innerHTML}, it never
interprets HTML --- safe by design.

\textbf{Transaction} --- A group of database operations that must
succeed together or fail together. Wrapped in
\texttt{beginTransaction()} / \texttt{commit()} /
\texttt{rollBack()}. In Bibliotheca, \texttt{Book::delete()} uses
a transaction: deleting the \texttt{book\_author} records and the
\texttt{book} record must be atomic --- if one fails, neither
should persist.

\textbf{URL rewriting} --- Transforming \texttt{/publishers} into
\texttt{index.php?route=publishers} so the user sees clean URLs
while the server uses query parameters. Done by Apache's
\texttt{mod\_rewrite} via \texttt{.htaccess}.

\textbf{Whitelist} --- A list of explicitly allowed values. The
router in \texttt{index.php} checks the route against a whitelist
of valid pages. Anything not on the list gets a 404. Safer than a
blacklist (blocking known bad values) because unknown values are
rejected by default.

\textbf{XSS (Cross-Site Scripting)} --- An attack where malicious
JavaScript is injected into a page. Prevented by escaping output:
\texttt{htmlspecialchars()} in PHP, \texttt{textContent} in
JavaScript. Never render user input as HTML.

\textbf{YAGNI (You Aren't Gonna Need It)} --- Don't build what
you don't need yet. Solve today's problem today.

\newpage

\hypertarget{bibliography}{%
\chapter{Bibliography}\label{bibliography}}

\hypertarget{the-terminal-1}{%
\section{The terminal}\label{the-terminal-1}}

\begin{itemize}
\tightlist
\item
  William Shotts, \emph{The Linux Command Line}, Sixth Internet
  Edition. Free at https://linuxcommand.org
\end{itemize}

\hypertarget{apache}{%
\section{Apache}\label{apache}}

\begin{itemize}
\tightlist
\item
  \emph{Apache .htaccess Tutorial} ---
  https://httpd.apache.org/docs/2.4/howto/htaccess.html
\end{itemize}

\hypertarget{php}{%
\section{PHP}\label{php}}

\begin{itemize}
\tightlist
\item
  \emph{PHP Manual} --- https://www.php.net/manual/en/
\item
  \emph{PHP: The Right Way} --- https://phptherightway.com
\item
  \emph{PHP Delusions} --- https://phpdelusions.net
\end{itemize}

\hypertarget{javascript-html-css}{%
\section{JavaScript, HTML, CSS}\label{javascript-html-css}}

\begin{itemize}
\tightlist
\item
  Marijn Haverbeke, \emph{Eloquent JavaScript}, 4th Edition. Free
  at https://eloquentjavascript.net
\item
  \emph{MDN Web Docs} --- https://developer.mozilla.org
\end{itemize}

\hypertarget{sql-and-sqlite}{%
\section{SQL and SQLite}\label{sql-and-sqlite}}

\begin{itemize}
\tightlist
\item
  \emph{SQLite Documentation} --- https://sqlite.org/docs.html
\item
  \emph{SQLite Tutorial} --- https://www.sqlitetutorial.net
\end{itemize}

\hypertarget{http}{%
\section{HTTP}\label{http}}

\begin{itemize}
\tightlist
\item
  \emph{MDN: HTTP Guide} ---
  https://developer.mozilla.org/en-US/docs/Web/HTTP
\end{itemize}

\hypertarget{git}{%
\section{Git}\label{git}}

\begin{itemize}
\tightlist
\item
  Scott Chacon, Ben Straub, \emph{Pro Git}, 2nd Edition. Free at
  https://git-scm.com/book
\end{itemize}

\hypertarget{the-classics}{%
\section{The classics}\label{the-classics}}

\begin{itemize}
\tightlist
\item
  Brian W. Kernighan, Dennis M. Ritchie, \emph{The C Programming
  Language}, 2nd Edition, Prentice Hall, 1988. You will not need
  C for web development, but this book will shape the way you
  think as a programmer.
\end{itemize}

\hypertarget{security-1}{%
\section{Security}\label{security-1}}

\begin{itemize}
\tightlist
\item
  \emph{OWASP Top Ten} --- https://owasp.org/www-project-top-ten/
\end{itemize}

\newpage

\hypertarget{the-apocryphal-chapter-beyond-the-launch}{%
\chapter{The Apocryphal Chapter --- Beyond the
Launch}\label{the-apocryphal-chapter-beyond-the-launch}}

\begin{quote}
\emph{And now for something completely different.} --- Monty
Python
\end{quote}

You made it through ten lessons. You built a web application from
scratch --- database, backend, API, frontend --- with nothing
between you and the code. No framework decided for you, no
library hid the mechanics.

Now what?

\hypertarget{what-you-actually-learned}{%
\section{What you actually
learned}\label{what-you-actually-learned}}

You did not just learn PHP, JavaScript, and SQL. You learned a
way of thinking:

\begin{itemize}
\tightlist
\item
  \textbf{A request is a conversation.} The browser asks, the
  server answers. HTTP methods are the vocabulary. JSON is the
  language. One URL, different verbs --- that is REST.
\item
  \textbf{Data flows in one direction.} JavaScript calls the
  controller, the controller calls the model, the model calls
  DBMS, DBMS talks to SQLite --- and back. Each layer has one
  job. If you respect the boundaries, the code stays simple.
\item
  \textbf{Security is not a feature.} It is a habit. Prepared
  statements, \texttt{textContent}, CSRF tokens, server-side
  validation --- these are not optional extras. They are how you
  write code. Each defense stands on its own.
\item
  \textbf{The database is the truth.} Everything else --- the
  controller, the frontend, the cache --- is a representation.
  When in doubt, ask the database.
\end{itemize}

\hypertarget{what-bibliotheca-does-not-do}{%
\section{What Bibliotheca does not
do}\label{what-bibliotheca-does-not-do}}

Bibliotheca has authentication, search, pagination, sorting, CSRF
protection, and input validation. That is more than most
tutorials cover. But a production application would still need:

\textbf{Multi-user roles.} We have a single admin. A real system
needs roles (admin, editor, viewer), permissions per resource,
and audit logs of who changed what.

\textbf{Server-side pagination.} Our pagination is client-side
--- all data is loaded at once, then sliced in JavaScript. With
ten thousand books, you need \texttt{LIMIT} and \texttt{OFFSET}
in SQL, and the API must accept \texttt{?page=3\&per\_page=20}.

\textbf{Full-text search.} Our search filters in memory. A real
search needs \texttt{FTS5} in SQLite (or Elasticsearch for
scale), ranking, highlighting, and fuzzy matching.

\textbf{Error logging and monitoring.} Our API returns
appropriate status codes (400, 401, 409, 422) with error
messages, but it does not log errors to a file or notify anyone.
A production app needs structured logging, alerts, and
dashboards.

\textbf{Deployment.} Moving from \texttt{localhost} to a real
server: domain names, HTTPS, firewall rules, backups, monitoring.
The code is the easy part.

\hypertarget{frameworks-when-and-why}{%
\section{Frameworks --- when and
why}\label{frameworks-when-and-why}}

At this point, someone will tell you: ``Just use Laravel'' or
``Just use React.'' And they are not wrong --- for production
code.

But now you know what a framework does for you:

\begin{itemize}
\tightlist
\item
  \textbf{A router} --- you built one in \texttt{index.php}. A
  whitelist, a parameter, a \texttt{require}. Laravel's router
  does the same thing with more features.
\item
  \textbf{An ORM} --- you built models with \texttt{insert},
  \texttt{getAll}, \texttt{update}, \texttt{delete}, and raw SQL.
  Eloquent writes the SQL for you. But when it generates a slow
  query, you will know why --- because you understand
  \texttt{JOIN} and \texttt{INDEX}.
\item
  \textbf{A template engine} --- you used \texttt{createElement}
  and \texttt{textContent} to build the DOM. React does the same
  with JSX, plus reactivity. But the DOM manipulation underneath
  is identical.
\item
  \textbf{Middleware} --- your controllers check the HTTP method
  with
  \texttt{\$\_SERVER{[}\textquotesingle{}REQUEST\_METHOD\textquotesingle{}{]}}.
  Express or Laravel route middleware does this declaratively.
\item
  \textbf{A REST client} --- you used \texttt{fetch} with
  \texttt{async/await}, built a \texttt{payload}, and chose the
  HTTP method. Axios or other libraries wrap the same mechanism
  with convenience methods.
\end{itemize}

A framework is not magic. It is someone else's
\texttt{index.php}, someone else's \texttt{DBMS.php}, someone
else's \texttt{Publisher.js} --- packaged, tested, and
documented. You can use it with confidence now, because you know
what it replaces.

The danger is using a framework \emph{without} understanding the
layer below. You can configure a router without knowing what
\texttt{mod\_rewrite} does. You can call an ORM without knowing
what a prepared statement is. But when something breaks --- and
it will --- you will stare at the error with no idea where to
look.

That is the difference between a developer and someone who uses
developer tools.

\hypertarget{the-age-of-ai}{%
\section{The age of AI}\label{the-age-of-ai}}

You are learning to program in an era where machines can write
code. That changes the job, not the need.

An AI can generate a CRUD endpoint in seconds. But it cannot
think for you. You are the one who defines the domain, decides
which entities exist, what rules govern them, and how they relate
to each other. You model the problem. The machine writes code ---
you decide what code should be written, and whether it is right.
If you give up that responsibility, you have no reason to be
there. That should be unsettling enough to keep you learning.
Remember this always.

The programmers who will thrive are not the ones who type
fastest. They are the ones who understand what the machine wrote,
can verify it is correct, and know when to throw it away and
start over.

\emph{Learn to read, write and compute --- before it's too late.}

\hypertarget{section}{%
\section{42}\label{section}}

The code you wrote in these lessons is yours. Read it again in
six months. You will understand it differently --- not because
the code changed, but because you did.

That is the answer.

\newpage

\end{document}
